\documentclass[11pt]{ctexart}
\usepackage[a4paper,margin=1in]{geometry}
\usepackage{graphicx}
\usepackage{xcolor}
\usepackage{hyperref}
\definecolor{refgray}{gray}{0.45}
\title{\textbf{Market Views｜全球投行叙事汇编}\\[0.4em]{\large 美联储框架重塑与人民币估值重估并行，AI硬件需求从GPU向CPU、内存与封装协同扩散，个股与指数波动率背离接近历史极值}}
\author{KC桌面 自动生成}
\date{260716}
\begin{document}
\maketitle
\tableofcontents
\newpage
\section{一页摘要}
\begin{itemize}
  \item 美联储主席Warsh以“防扩散”原则重塑通胀治理框架，缩表节奏审慎但非转向宽松，市场部分参与者误读为宽松信号。
  \item 人民币中间价模型投影大幅下移，多模型交叉验证显示人民币对欧元仍被低估约8\%-20\%，政策层对当前汇率水平的接受度在提高。
  \item 个股已实现波动率显著高于指数波动率，背离幅度接近历史极值，半导体板块是主要推手，BofA认为相关性回升将导致指数波动率快速扩大，而MS认为调整是暂时的。
  \item AI硬件需求正从单一的GPU算力驱动，转向CPU、内存、先进封装与光互连的协同增长，代理式AI架构将重塑计算重心，CPU增量机会达2380亿美元但市场尚未充分定价。
  \item 台积电与ASML财报共同指向AI需求从芯片制造端向设备端传导，先进制程产能利用率超载，供需缺口至少延续至2028年，毛利率上行成为核心边际变化。
  \item 中国互联网公司普遍从追求规模转向盈利优先，但抖音与美团在本地生活的竞争格局分化，抖音核销率（57-58\%）远低于美团（>80\%），净交易额差距更大。
  \item 全球汽车市场区域分化显著：中国纯电渗透率创新高但国内需求疲弱，欧洲受政策与平价车型驱动增长强劲，美国电动化进程遇阻，中国品牌海外扩张成为结构性亮点。
  \item GLP-1市场2031年将达千亿美元，礼来占据约70\%份额，口服药并非颠覆者；心血管领域Lp(a)靶点面临关键验证，HORIZON试验若成功将重塑赛道估值。
\end{itemize}
\section{全球投行叙事汇编}
今日纳入 87 篇投行/券商研究，按 19 个市场、资产与产业主题系统整理。
\newpage
\subsection{宏观与利率：美联储框架重塑与人民币估值重估}
\textbf{美联储主席Warsh正以“防扩散”原则重塑通胀治理框架，缩表节奏将更审慎但非转向宽松；同时，人民币中间价模型投影大幅下移，多模型交叉验证显示人民币对欧元仍被低估约15\%，但政策层对当前汇率水平的接受度在提高。}
\subsubsection*{主流共识}
\begin{itemize}
  \item 美联储主席Warsh强调通胀治理的‘防扩散’原则，对单月CPI数据持谨慎态度，拒绝过度解读。
  \item 人民币中间价模型投影值大幅下调，显示定价基准正在主动下移，逆周期因子收敛。
  \item 多模型交叉验证显示人民币对欧元仍被低估约8\%-20\%，外部均衡模型指向结构性矛盾。
  \item 中国‘十五五’扩大消费规划以供给侧升级为核心，短期提振效果有限。
\end{itemize}
\subsubsection*{分歧与条件差异}
\begin{itemize}
  \item 市场部分参与者将Warsh对缩表的审慎态度解读为宽松信号，但UBS认为这更接近‘操作纪律升级’，而非方向性转向。
  \item 人民币估值分歧：DB的十种模型显示低估8\%-20\%，但NOM模型投影显示中间价可能更快向市场供需靠拢，政策接受度提高。
  \item 日本CPI基期修订对通胀读数的影响方向存在分歧：短期上修但中期可能形成向下牵引。
\end{itemize}
\subsubsection*{机构观点}
\paragraph{NOM} NOM的美元/人民币中间价模型投影值从6.7990下调至6.7681，下调309个基点，计入逆周期因子后为6.7745，较前一日实际中间价低245个基点。模型投影已明显低于现货收盘价，表明中间价定价基准正在主动下移，逆周期因子收敛至仅64个基点，显示政策层对当前汇率水平的接受度提高。
{\small\color{refgray}数据：模型投影值6.7681，较前次下调309个基点\par}
{\small\color{refgray}数据：计入逆周期因子后投影6.7745，较实际中间价低245个基点\par}
{\small\color{refgray}数据：逆周期因子仅64个基点，远低于历史极端时期\par}
{\small\color{refgray}边际变化：模型投影从高于现货收盘价转为低于现货收盘价，定价基准发生位移，逆周期因子收敛至历史低位。\par}
\paragraph{UBS} UBS认为美联储主席Warsh在国会听证会上释放了三个关键信号：缩表节奏将比市场预期更审慎，强调‘深思熟虑、公开宣布、逐步推进’；明确押注AI驱动生产率提升，将其写入货币政策因果链；用‘任务组’替代‘点阵图’重塑通胀治理框架。美联储正从‘数据依赖’转向‘框架依赖’，市场对此定价远未充分。
{\small\color{refgray}数据：Warsh明确否定回到2006年小规模资产负债表\par}
{\small\color{refgray}数据：缩表调整需‘公开宣布’且‘相当长时间落地’\par}
{\small\color{refgray}数据：AI被视为类似90年代互联网的生产力爆发机会\par}
{\small\color{refgray}边际变化：市场此前预期Warsh延续前任缩表节奏，但听证会显示其操作空间主动收窄，且强调框架依赖而非数据依赖。\par}
\paragraph{DB} DB通过十种模型交叉验证欧元/人民币合理估值，结论一致：人民币对欧元仍被低估约8\%-20\%。巨无霸指数暗示低估46\%，OECD购买力平价约34\%，去趋势调整后约20\%，相对购买力平价（PPI口径）约14\%，DBeer模型约8\%。外部均衡模型显示中国顺差扩大与欧元区逆差的结构性矛盾仍在调整。
{\small\color{refgray}数据：巨无霸指数暗示人民币对欧元低估约46\%\par}
{\small\color{refgray}数据：OECD购买力平价暗示低估约34\%\par}
{\small\color{refgray}数据：去趋势调整后低估约20\%\par}
{\small\color{refgray}数据：DBeer模型给出约8\%低估\par}
{\small\color{refgray}边际变化：此前市场关注单一指标偏差，DB首次用十种模型交叉验证，形成‘模型共识区间’，且低估幅度接近2005-2008年极端水平。\par}
\paragraph{GS} GS分析中国‘十五五’扩大消费规划，核心是供给侧升级而非需求侧刺激。规划设定了到2030年社零总额约60万亿元的目标，年均名义增速3.7\%，低于‘十四五’的5.0\%。重点包括服务消费升级、耐用品置换需求、新消费业态培育，但短期提振效果有限，因劳动力市场调整、收入预期改善和房产财富效应修复需要时间。
{\small\color{refgray}数据：1-5月社零同比仅增1.4\%，远低于去年同期的5\%\par}
{\small\color{refgray}数据：2030年社零目标60万亿元，年均名义增速3.7\%\par}
{\small\color{refgray}数据：服务消费占居民消费比重仍远低于主要发达地区\par}
{\small\color{refgray}边际变化：政策从‘补贴驱动’转向‘供给升级驱动’，不再依赖一次性补贴，而是通过提升供给质量释放潜在需求。\par}
\paragraph{GS} GS测算日本CPI基期修订将上修2026年上半年核心CPI和新核心CPI读数，主因权重调整：高中学费、保育费权重下降使免费教育政策对CPI的负向影响缩小；大米价格权重从0.62\%升至1.01\%，将尤其抬升2026年1-3月读数。但中期看，新基期可能对通胀预测形成向下牵引，因大米价格后续回落且手机费用权重降低。
{\small\color{refgray}数据：大米价格权重从0.62\%升至1.01\%\par}
{\small\color{refgray}数据：手机费用权重从0.90\%降至0.58\%\par}
{\small\color{refgray}数据：基期修订将于8月7日启动\par}
{\small\color{refgray}边际变化：市场此前关注基期修订对通胀读数的上修影响，GS指出中期可能形成向下牵引，改变通胀叙事节奏。\par}
\paragraph{GS} GS分析美联储主席Warsh在国会听证会上的表态，核心是‘防扩散’原则：确保短期价格变化不会演变为普遍的价格水平变动。Warsh明确表示6月CPI走软只是‘一个数据点’，拒绝过度解读，并强调对持续高通胀‘零容忍’。缩表规模和持续时间值得重新评估，但调整将审慎推进。
{\small\color{refgray}数据：Warsh称6月CPI走软‘只是一个数据点’\par}
{\small\color{refgray}数据：强调‘零容忍’对持续高通胀\par}
{\small\color{refgray}数据：缩表‘规模和持续时间’均值得重新评估\par}
{\small\color{refgray}边际变化：市场此前预期Warsh可能因通胀数据改善而软化立场，但听证会显示其通胀治理逻辑更严格，且拒绝将单月数据作为政策转向依据。\par}
\paragraph{DB} 人民币国际化正在经历一个关键转折。近期研究指出，2026年上半年熊猫债发行量突破1600亿元人民币，同比增长超过60\%，而离岸点心债市场的外国发行人占比已超过25\%。这两个市场正在形成截然不同但又相互补充的“双引擎”结构——熊猫债正从境内融资工具转变为全球人民币资金来源，点心债则加速成为国际参与者的核心资产池。这一变化的背后，是相关规划将人民币国际化明确列为...
{\small\color{refgray}数据：熊猫债市场的增长逻辑已经发生根本性变化。报告数据显示，2026年上半年超过一半的募集资金被用于离岸用途，这一比例从过去的不到40\%跃升至53\%。这意味着熊猫债不再只是境外机构进入境内市场的通道，而是变成了真正可全球流通的人民币融资工具。\par}
{\small\color{refgray}数据：但市场结构仍存在明显短板。几乎所有发行期限都在3年以内，这并非监管限制，而是境内参与者结构决定的——银行主导的短期资金偏好使得长期限债券缺乏需求。报告认为，熊猫债的战略潜力仍被短期化特征所约束，资金类发行人占比过高也限制了发行主体的多元化。\par}
\subsubsection*{关键数据}
\begin{itemize}
  \item NOM模型投影值6.7681，较前次下调309个基点
  \item 逆周期因子仅64个基点，远低于历史极端时期
  \item DB十种模型显示人民币对欧元低估8\%-20\%
  \item 中国1-5月社零同比仅增1.4\%
  \item 十五五社零目标年均增速3.7\%，低于十四五的5.0\%
  \item 日本大米价格权重从0.62\%升至1.01\%
  \item Warsh称6月CPI走软‘只是一个数据点’
  \item Warsh强调对持续高通胀‘零容忍’
\end{itemize}
\begin{figure}[htbp]
\centering
\includegraphics[width=0.88\linewidth]{figures/fig_082.jpg}
\caption{Figure 2 - Figure 2: Euro Area trade deficit with China in goods is widening}
\end{figure}
\begin{figure}[htbp]
\centering
\includegraphics[width=0.88\linewidth]{figures/fig_113.jpg}
\caption{Exhibit 2 - Exhibit 1: Core and New Core CPI Inflation for January-June 2026 Could Be Revised Up at the August CPI Base Year Revision Exhibit 2: Weights of CPI Main Categories 2025 New Basis vs. 2020 Current Basis (per 10000)}
\end{figure}
\begin{figure}[htbp]
\centering
\includegraphics[width=0.88\linewidth]{figures/fig_083.jpg}
\caption{Figure 3 - Figure 3: The relative price of a burger suggests the yuan is close to 50\% undervalued against the euro}
\end{figure}
\begin{figure}[htbp]
\centering
\includegraphics[width=0.88\linewidth]{figures/fig_084.jpg}
\caption{Figure 4 - Figure 5: CNY had trended less cheap over time on Big Mac Index...}
\end{figure}
{\small\color{refgray}本节综合 7 篇研究；涉及 NOM、UBS、DB、GS。}
\newpage
\subsection{权益与波动率：个股分化加剧，指数波动被低估}
\textbf{当前市场呈现罕见的个股与指数波动率背离，半导体板块是核心驱动；AI行情进入季节性调整但非趋势逆转，盈利扩散支撑市场领导权从少数龙头向外扩展。}
\subsubsection*{主流共识}
\begin{itemize}
  \item 个股已实现波动率显著高于指数波动率，背离幅度接近历史极值，半导体板块是主要推手。
  \item AI行情进入夏季调整，主要受季节性仓位调整和外资获利了结影响，而非基本面逆转。
  \item 盈利修正广度维持高位，标普500中位数公司实现双位数EPS增长，盈利扩散正在加速。
\end{itemize}
\subsubsection*{分歧与条件差异}
\begin{itemize}
  \item 对指数波动率上行风险的判断：BofA认为相关性回升将导致指数波动率快速扩大，而MS认为调整是暂时的，趋势未结束。
  \item 对AI行情调整深度的分歧：MS认为调整窗口约6周，9月初可能企稳；但BofA指出公司波动率已升至科网泡沫水平，若估值继续扩张，背离可能进一步扩大。
  \item 对欧洲周期股对冲价值的看法：BofA认为SX5E与斯托克600的波动率价差处于历史低位，对冲成本便宜；MS未涉及欧洲市场。
\end{itemize}
\subsubsection*{机构观点}
\paragraph{BofA} 美国市场公司已实现波动率升至1990年代末科网泡沫后期水平，但指数波动率相对温和，两者背离处于历史极端区间。半导体板块（如MU、AMD）的不稳定性推高公司波动率，同时半导体与其他科技股相关性降至接近历史最低，压制指数波动率。若相关性回升，指数波动率面临显著上行风险，尤其考虑到季节性挑战期。欧洲斯托克50指数与斯托克600的隐含波动率价差收窄至历史低位，低于已实现波动率价差，通过看跌期权互换对冲欧洲周期板块风险的成本异常低廉。
{\small\color{refgray}数据：公司已实现波动率升至1990年代末科网泡沫后期水平\par}
{\small\color{refgray}数据：半导体与其他大型科技股相关性降至接近历史最低\par}
{\small\color{refgray}数据：SX5E与SX6E隐含波动率价差低于已实现波动率价差\par}
{\small\color{refgray}数据：半导体BRI（泡沫不确定性指标）走高\par}
{\small\color{refgray}边际变化：从关注指数波动率转向个股与指数波动率裂口，强调半导体脱钩对市场结构的深远影响，并提示欧洲周期股对冲成本处于历史低位。\par}
\paragraph{MS} 市场领导权扩散有盈利数据支撑：等权标普500自5月中旬跑赢市值加权，标普1500中位数公司实现双位数EPS增长，为后疫情复苏以来最强。盈利修正广度在7月季节性回落中仍维持+24\%高位，可选消费和交通运输板块相对优势明显。AI采用率加速，40\%的AI采用公司报告可量化收益（一年前21\%），受益者从半导体扩散至更广泛行业。日本AI相关公司回调属季节性仓位调整和外资获利了结，外资上半年净买入10万亿日元后回撤正常，历史数据显示8-9月外资通常净流出，预计9月初调整结束。
{\small\color{refgray}数据：标普1500中位数公司EPS增速双位数，为后疫情复苏以来最强\par}
{\small\color{refgray}数据：盈利修正广度+24\%，7月季节性回落中仍处历史高位\par}
{\small\color{refgray}数据：40\% AI采用公司报告可量化收益，一年前仅21\%\par}
{\small\color{refgray}数据：外资上半年净买入日本现货10万亿日元\par}
{\small\color{refgray}边际变化：从聚焦AI龙头转向盈利扩散和AI采用率加速，强调日本AI调整是季节性而非趋势逆转，并提示外资调仓可能提前至7月。\par}
\subsubsection*{关键数据}
\begin{itemize}
  \item 公司已实现波动率升至1990年代末科网泡沫后期水平
  \item 半导体与其他大型科技股相关性降至接近历史最低
  \item 标普1500中位数公司EPS增速双位数
  \item 盈利修正广度+24\%
  \item 40\% AI采用公司报告可量化收益（一年前21\%）
  \item 外资上半年净买入日本现货10万亿日元
  \item SX5E与SX6E隐含波动率价差低于已实现波动率价差
\end{itemize}
\begin{figure}[htbp]
\centering
\includegraphics[width=0.88\linewidth]{figures/fig_015.jpg}
\caption{Exhibit 2 - Exhibit 2: Latest\textbackslash{}* stress across GFSI sub-components Tibor-OIS is the most stressed while sub-IG foreign sovereign bond spreads are the least stressed \#\# Exhibit 3: Change\textbackslash{}*\textbackslash{}* in stress across GFSI sub-components Euribor-OIS was}
\end{figure}
\begin{figure}[htbp]
\centering
\includegraphics[width=0.88\linewidth]{figures/fig_054.jpg}
\caption{Exhibit 1 - Exhibit 1: 2Q Reporting Season by Market Cap Note: Includes 451 of S\&P 500 companies and 95\% of S\&P market cap reporting from July 10 to Aug. 30 Exhibit 2: 2Q Reporting Season by Company Count Note: Includes 451 of S\&P 500 comp}
\end{figure}
\begin{figure}[htbp]
\centering
\includegraphics[width=0.88\linewidth]{figures/fig_060.jpg}
\caption{Exhibit 1 - Exhibit 3: Average Monthly TOPIX Returns: August Tends to Deliver Negative Returns}
\end{figure}
\begin{figure}[htbp]
\centering
\includegraphics[width=0.88\linewidth]{figures/fig_062.jpg}
\caption{Exhibit 2 - Exhibit 2: Percentage of Months with Positive Net Purchases by Foreign Investors, by Investor Type Exhibit 3: Average Monthly TOPIX Returns: August Tends to Deliver Negative Returns Exhibit 4: By Investor Type: Foreign Investor}
\end{figure}
{\small\color{refgray}本节综合 3 篇研究；涉及 BofA、MS。}
\newpage
\subsection{AI硬件需求结构分化：从GPU算力到CPU、内存与封装协同}
\textbf{AI硬件需求正从单一的GPU算力驱动，转向CPU、内存、先进封装与光互连的协同增长，代理式AI架构将重塑计算重心，供应链份额与产能分配进入再平衡阶段。}
\subsubsection*{主流共识}
\begin{itemize}
  \item AI基础设施资本开支持续加速，前五大云厂商2026年资本开支预计同比增长85\%，高于2025年的72\%，数据中心交换机市场保持50\%以上同比增长。
  \item 先进制程产能持续紧张，台积电2027年资本支出可能上调至802亿美元以上，ASML计划2028年EUV产能较2027年再增30\%，设备端订单动能强劲。
  \item 存储周期正在拉长而非急剧下行，DRAM合约价预计2026年四季度见顶，但市场定价理性，LTA未被重新定价。
\end{itemize}
\subsubsection*{分歧与条件差异}
\begin{itemize}
  \item 代理式AI是否将导致计算瓶颈从GPU向CPU和内存迁移，MS认为CPU增量机会达2380亿美元，但市场尚未充分定价。
  \item CPO交换机放量时间表存在分歧，BofA预计英伟达CPO交换机2028年才大规模爬坡，而部分供应链数据已显示提前迹象。
  \item 日本半导体设备与材料公司的利润增速（EPS CAGR 20.3\%）显著快于营收（13.2\%），但市场对其定价权可持续性看法不一。
  \item AI服务器HDI板份额格局正在调整，欣兴电子有望从不到20\%升至40\%，但胜宏科技的技术问题是否仅为暂时性仍存争议。
\end{itemize}
\subsubsection*{机构观点}
\paragraph{BofA} 全球交换机市场1Q26同比增长36.1\%，数据中心占66\%，AI后端交换机暴增121\%占数据中心需求的38\%。Cisco在AI后端从零起步拿下7.9\%份额，主要来自Meta 800G叶交换机。园区交换机同比增15.1\%，Cisco以53.5\%份额主导，Catalyst 4000/6000停售将推动换机需求。前五大云厂商2026年资本开支预计同比增85\%，高于2025年的72\%，前四大云厂商占数据中心交换机41\%份额，同比增90.4\%。
{\small\color{refgray}数据：全球交换机市场1Q26同比增36.1\%\par}
{\small\color{refgray}数据：AI后端交换机同比增121\%，占数据中心38\%\par}
{\small\color{refgray}数据：前五大云厂商2026年资本开支同比增85\%\par}
{\small\color{refgray}数据：Cisco在AI后端交换机份额从0升至7.9\%\par}
{\small\color{refgray}边际变化：AI后端交换机占比从一年前的26\%升至38\%，Cisco首次进入该市场并快速渗透，园区交换机恢复至9\%以上增长（历史常态1-2\%）。\par}
\paragraph{BofA} 光纤阵列单元（FAU）是CPO大规模量产的真正瓶颈，对准精度需控制在±0.3微米以下，比手机摄像头模组高一个数量级。英伟达CPO交换机2026/2027/2028年出货量预计为1.4万、5.2万和19万台，大规模爬坡在2028年后。FAU单台交换机物料成本占比5-10\%，2030年相关市场可达80亿美元。大立光在认证进度上领先其他新进入者。
{\small\color{refgray}数据：FAU对准精度要求±0.3微米，手机摄像头为±5微米\par}
{\small\color{refgray}数据：英伟达CPO交换机2028年出货19万台\par}
{\small\color{refgray}数据：FAU占CPO交换机BOM 5-10\%\par}
{\small\color{refgray}数据：2030年FAU市场TAM达80亿美元\par}
{\small\color{refgray}边际变化：市场此前低估FAU的进入壁垒，BofA明确指出微米级对准和自动化能力是核心瓶颈，而非光学引擎集成度。\par}
\paragraph{MS} 存储周期正在拉长，DRAM合约价预计4Q26见顶，但市场定价理性，LTA未被重新定价。AI资本支出仍处上升通道，但2Q26将是验证节点。代理式AI正从生成式过渡，计算瓶颈从GPU向CPU和内存迁移，到2030年可能创造2380亿美元CPU增量机会和221EB DRAM需求。台积电法说会三种情景中，AI需求可持续性是核心变量，最乐观情景下五年AI营收CAGR可上调至80\%。
{\small\color{refgray}数据：DRAM合约价预计4Q26见顶\par}
{\small\color{refgray}数据：代理式AI到2030年CPU增量机会2380亿美元\par}
{\small\color{refgray}数据：代理式AI拉动221EB DRAM需求\par}
{\small\color{refgray}数据：台积电最乐观情景AI五年CAGR上调至80\%\par}
{\small\color{refgray}边际变化：MS首次系统提出代理式AI将改变计算重心，CPU和内存增量需求被市场低估，台积电法说会关注点从季度营收转向AI需求时间跨度。\par}
\paragraph{NOM} 5月日本封装基板出货金额同比增36\%至278亿日元，单价同比升23\%至每平方米135.6万日元，均创历史新高。先进封装技术重心从互联密度转向供电与散热，Intel EMIB-T在硅桥中加入TSV和MIM电容，电压降改善68-80\%，支持四芯片、120x120mm基板、最多12颗HBM4。台积电通过3DFabric联盟在供电和散热上建立优势。
{\small\color{refgray}数据：5月封装基板出货金额同比增36\%至278亿日元\par}
{\small\color{refgray}数据：单价同比升23\%至每平方米135.6万日元\par}
{\small\color{refgray}数据：EMIB-T电压降改善68-80\%\par}
{\small\color{refgray}数据：支持四芯片、120x120mm基板、最多12颗HBM4\par}
{\small\color{refgray}边际变化：封装基板单价持续走高反映高价值产品占比扩大，技术论文揭示供电瓶颈成为先进封装新焦点，EMIB-T方案有望2026年量产。\par}
\paragraph{UBS} 7月初路演反馈显示，市场对电子元器件板块兴趣仍高，但开始区分谁能真正在AI数据中心拿到订单。多数读者看好MLCC基本面，但担忧估值偏高和报价波动加大，村田优于太阳诱电因其AI中期增长可见度更高。罗姆评级上调后，市场关注其内在价值（剔除东芝持股后核心业务被低估）及Kioxia关联。
{\small\color{refgray}数据：28场路演交流于香港和新加坡\par}
{\small\color{refgray}数据：村田AI中期增长可见度高于太阳诱电\par}
{\small\color{refgray}数据：罗姆核心业务剔除东芝持股后仍被低估\par}
{\small\color{refgray}数据：MLCC基本面共识强但估值担忧增加\par}
{\small\color{refgray}边际变化：市场从全面看多转向区分个股阿尔法，村田与太阳诱电的估值差异反映AI业务可见度定价，罗姆评级上调后关注点从时机转向内在价值。\par}
\paragraph{HSBC} 台积电2027年资本支出预估上调至802亿美元，远高于市场共识的652亿美元，但可能仍不足以解决先进制程产能紧张。敏感性分析显示，若资本支出增至850-1000亿美元，2029年营收可在25\% CAGR基础上再增4-19\%。2Q26毛利率预计达68\%高于指引上沿，但3Q26因海外厂和2nm爬坡可能回落至67\%。
{\small\color{refgray}数据：2027年资本支出预估802亿美元，市场共识652亿\par}
{\small\color{refgray}数据：2Q26毛利率预计68\%，高于指引上沿65.5-67.5\%\par}
{\small\color{refgray}数据：3Q26毛利率可能回落至67\%\par}
{\small\color{refgray}数据：若资本支出850-1000亿，2029年营收额外增4-19\%\par}
{\small\color{refgray}边际变化：市场焦点从短期毛利率转向长期产能扩张，HSBC认为802亿美元资本支出可能仍不够，扩产规模成为增长上限的关键变量。\par}
\paragraph{JPM} ASML 2Q26营收93.3亿欧元超指引，毛利率54\%高于51-52\%指引。2026年全年收入指引从360-400亿欧元上调至430-450亿欧元，中位数440亿欧元较此前高16\%。2028年EUV产能计划达约110台，浸没式DUV约210台，均较2027年再增30\%，基于此2028年EPS可能超65欧元。半导体设备与器件走向分化，设备端订单动能加速且客户结构多元化，器件端面临存储涨价成本压力。
{\small\color{refgray}数据：ASML 2Q26营收93.3亿欧元，超84-90亿指引\par}
{\small\color{refgray}数据：2026年收入指引上调至430-450亿欧元\par}
{\small\color{refgray}数据：2028年EUV产能约110台，较2027年增30\%\par}
{\small\color{refgray}数据：2028年EPS可能超65欧元\par}
{\small\color{refgray}边际变化：ASML 2028年产能指引高于JPM自身预期，市场此前对长期增长潜力存在系统性低估；设备与器件分化加剧，设备端估值仍处十年中位数。\par}
\paragraph{JPM} 日本半导体设备与材料公司利润增速显著快于营收，平均EPS三年CAGR 20.3\% vs 营收CAGR 13.2\%，差距近7个百分点，反映定价权向头部集中。东京电子EPS CAGR 27.7\%高于营收23.2\%，Advantest EPS CAGR 26.7\%高于营收23.1\%。技术材料类公司平均营收增速仅6.8\%，EPS增速13.1\%，利润改善幅度更小。
{\small\color{refgray}数据：日本半导体设备与材料平均EPS CAGR 20.3\%\par}
{\small\color{refgray}数据：平均营收CAGR 13.2\%\par}
{\small\color{refgray}数据：东京电子EPS CAGR 27.7\% vs 营收23.2\%\par}
{\small\color{refgray}数据：Advantest EPS CAGR 26.7\% vs 营收23.1\%\par}
{\small\color{refgray}边际变化：利润增速跑赢营收7个百分点，商业模式从设备出货量驱动转向工艺know-how和客户粘性驱动，技术壁垒货币化能力增强。\par}
\paragraph{JPM} 下一代AI服务器运算托盘板供应商格局调整，欣兴电子有望拿下约40\%份额，较当前世代初期不到20\%大幅提升。胜宏科技因技术细节问题份额暂时回落，但问题不严重，不会推迟整体GPU服务器项目。欣兴产能本就紧张，因资源向光收发模块mSAP板等新领域倾斜，AI订单增量与光通信等非AI领域争夺高端PCB产能。
{\small\color{refgray}数据：欣兴电子下一代AI服务器运算托盘板份额约40\%\par}
{\small\color{refgray}数据：当前世代初期份额不到20\%\par}
{\small\color{refgray}数据：胜宏科技因技术细节问题份额暂时回落\par}
{\small\color{refgray}数据：欣兴产能紧张因资源向光模块mSAP板倾斜\par}
{\small\color{refgray}边际变化：AI服务器HDI板从一家独大转向多家分散，份额调整基于产能调配和短期技术问题，非技术路线根本性颠覆，但供应链议价结构正在变化。\par}
\subsubsection*{关键数据}
\begin{itemize}
  \item 全球交换机市场1Q26同比增36.1\%，AI后端交换机同比增121\%
  \item 前五大云厂商2026年资本开支同比增85\%，高于2025年的72\%
  \item FAU对准精度要求±0.3微米，CPO交换机2030年TAM达80亿美元
  \item 代理式AI到2030年CPU增量机会2380亿美元，DRAM需求221EB
  \item 台积电2027年资本支出预估802亿美元，远高于市场共识652亿
  \item ASML 2028年EUV产能约110台，较2027年增30\%
  \item 日本半导体设备与材料平均EPS CAGR 20.3\% vs 营收CAGR 13.2\%
  \item 欣兴电子AI服务器运算托盘板份额从不到20\%升至约40\%
  \item 5月日本封装基板出货金额同比增36\%至278亿日元，单价创历史新高
  \item DRAM合约价预计4Q26见顶，但周期正在拉长
\end{itemize}
\begin{figure}[htbp]
\centering
\includegraphics[width=0.88\linewidth]{figures/fig_028.jpg}
\caption{Exhibit 1 - Exhibit 1: We expect Nvidia CPO switch shipment to reach 14k/52k/190k units in 2026/27/28E Nvidia CPO switch shipment, 2026-30E BofA GLOBAL RESEARCH Exhibit 3: We estimate FAU value per switch could reach US\textbackslash{}\$6-7k Nvidia CPO swit}
\end{figure}
\begin{figure}[htbp]
\centering
\includegraphics[width=0.88\linewidth]{figures/fig_065.jpg}
\caption{Exhibit 1 - \textbackslash{}- Scenario 3:TSMC guides for 3Q26 revenue to be up 5\%-10\% Q/Q in USD, full-year revenue growth updated to \textbackslash{}\textasciitilde{}35\% Y/Y, 3Q gross margin to be \textbackslash{}\textasciitilde{}66\%-67\% (down Q/Q). 2026 full-year capex guidance unchanged at US\textbackslash{}\$54-56bn, while five-year AI semi revenue CAGR is al}
\end{figure}
\begin{figure}[htbp]
\centering
\includegraphics[width=0.88\linewidth]{figures/fig_142.jpg}
\caption{Exhibit 2 - Exhibit 2: TSMC capex (USDm) estimate comparisons Exhibit 3: TSMC – capacity added per EUV equipment}
\end{figure}
\begin{figure}[htbp]
\centering
\includegraphics[width=0.88\linewidth]{figures/fig_048.jpg}
\caption{Exhibit 2 - Exhibit 2: We see 92 and 104 EUV Shipments for FY27 and FY28 respectively. \#\# 2Q26 Analog Playbook Exhibit 1: Cycle recovery is in play... Exhibit 2: ... with margin expansion on the horizon. The Americas}
\end{figure}
{\small\color{refgray}本节综合 12 篇研究；涉及 BofA、MS、NOM、UBS、HSBC、JPM。}
\newpage
\subsection{AI需求驱动先进制程与设备双轮景气}
\textbf{台积电与ASML的财报共同指向AI需求正在从芯片制造端向设备端传导，先进制程产能利用率超载、设备订单加速，供需缺口至少延续至2028年，毛利率上行成为核心边际变化。}
\subsubsection*{主流共识}
\begin{itemize}
  \item AI数据中心需求是台积电和ASML业绩超预期的共同核心驱动，且需求持续性已通过客户长期协议（LTA）得到确认。
  \item 先进制程（N3/N5）和High NA EUV设备产能利用率均处于满载甚至超负荷状态，供需缺口至少延续到2027-2028年。
  \item 毛利率上行趋势明确：台积电Q2毛利率逼近70\%，ASML Q3毛利率指引55-57\%，均显著超出市场预期，反映定价权和产品结构优化。
\end{itemize}
\subsubsection*{分歧与条件差异}
\begin{itemize}
  \item 台积电毛利率能否持续维持在67\%以上存在分歧：乐观方认为利用率溢价和效率改善可抵消N2及海外产能稀释；谨慎方担忧产能扩张带来的折旧压力。
  \item ASML产能扩张节奏（2026年65套Low NA EUV，2028年110套）是否足以满足需求：部分观点认为扩产速度仍慢于需求增长，另一部分认为客户资本开支承诺存在执行风险。
  \item CPU（NVIDIA Vera/AMD Venice）作为AI之外的新增长引擎，其贡献规模和时间点存在不确定性。
\end{itemize}
\subsubsection*{机构观点}
\paragraph{JPM} 台积电Q2营收1.27万亿新台币（环比+12\%，同比+36\%）落在美元指引高端，但更关键的是毛利率预计达69.5\%，显著超出65.5-67.5\%的官方指引区间。N3/N5利用率超120\%，热运行/超热运行订单溢价50\%/100\%，供需缺口至少延续至2028年。2026-28年资本支出预期上调至580/780/840亿美元，N3月产能2028年底达24万片，但缺口仍存。2027年预计N3/N2涨价8-10\%。CPU（NVIDIA Vera/AMD Venice）开始采用CoWoS封装，成为AI加速器之外的新增长引擎。
{\small\color{refgray}数据：台积电Q2营收1.27万亿新台币，环比+12\%，同比+36\%\par}
{\small\color{refgray}数据：Q2毛利率预计69.5\%，指引区间65.5-67.5\%\par}
{\small\color{refgray}数据：N3/N5利用率超120\%，热运行订单溢价50\%，超热运行溢价100\%\par}
{\small\color{refgray}数据：2026-28年资本支出预期580/780/840亿美元\par}
{\small\color{refgray}边际变化：毛利率预期从指引上限（67.5\%）上修至69.5\%，反映利用率超载和订单溢价超预期；CPU加入CoWoS封装需求，拓宽增长基础。\par}
\paragraph{MS} ASML Q2营收93亿欧元（超预期88亿），毛利率54\%，Q3指引营收115亿欧元（超预期104亿），毛利率55-57\%（超预期约52\%），2026全年指引440亿欧元（同比+34\%，超预期396亿）。AI需求是核心驱动，客户加速扩产并签署长期协议（LTA）。Low NA EUV产能规划：2026年约65套，2027年约85套，2028年约110套。High NA EUV进展超预期，英特尔18A制程已开始使用。毛利率改善源于高销量和产品组合优化（含IBM业务），2026全年毛利率指引约55\%，比市场共识高约260个基点。
{\small\color{refgray}数据：ASML Q2营收93亿欧元，超预期88亿\par}
{\small\color{refgray}数据：Q3营收指引115亿欧元，超预期104亿\par}
{\small\color{refgray}数据：Q3毛利率指引55-57\%，市场预期约52\%\par}
{\small\color{refgray}数据：2026全年营收指引440亿欧元，同比+34\%，市场预期396亿\par}
{\small\color{refgray}边际变化：Q3营收指引较市场预期高出11\%，毛利率指引较预期高约260个基点，反映AI需求对设备销售的拉动超预期；客户LTA签署确认需求持续性。\par}
\subsubsection*{关键数据}
\begin{itemize}
  \item 台积电Q2毛利率预计69.5\%，指引区间65.5-67.5\%
  \item N3/N5利用率超120\%，订单溢价50-100\%
  \item 台积电2026-28年资本支出580/780/840亿美元
  \item N3供需缺口约60万片，2028年底月产能24万片
  \item ASML Q3营收指引115亿欧元，超预期11\%
  \item ASML Q3毛利率指引55-57\%，超预期约260个基点
  \item ASML 2026全年营收指引440亿欧元，同比+34\%
  \item Low NA EUV产能：2026年65套，2028年110套
\end{itemize}
{\small\color{refgray}本节综合 2 篇研究；涉及 JPM、MS。}
\newpage
\subsection{AI与数据中心：融资结构分化与效率拐点}
\textbf{AI基础设施融资正从规模扩张转向结构质量，市场对回报周期的耐心在下降；同时，AI驱动的开发效率提升开始实质性改变IT服务成本结构，而欧洲电气化投资对电价的冲击被高估。}
\subsubsection*{主流共识}
\begin{itemize}
  \item AI相关债务发行规模已达3360亿美元，但市场情绪从全面接纳转向有选择地关注，偏好短期、有摊销机制的结构。
  \item 编码能力是模型智能迭代的核心飞轮，更强的编码能力正反哺模型研发，形成递归自我改进。
  \item AI基础设施投资焦点从GPU扩展至冷却、供电等配套，三菱重工等工业集团受益。
\end{itemize}
\subsubsection*{分歧与条件差异}
\begin{itemize}
  \item 数据中心债券的资产不确定性（续租、技术过时）是否被充分定价：不摊销结构在利率高位放大风险，但部分有施工保证的债券仍获认可。
  \item 欧洲电气化投资对电价的影响：市场普遍担忧推高电价，但GS认为35-40\%投资非通胀性，且需求增长摊薄固定成本。
  \item 中国AI模型商业化路径：UBS强调token ROI和编码agent驱动的货币化，但模型层盈利时间表仍不确定。
\end{itemize}
\subsubsection*{机构观点}
\paragraph{MS} AI债务发行已达3360亿美元，但市场正从接纳转向谨慎。超大规模企业发行激增、甲骨文评级下调后，AI债券利差走阔24bp。投资者偏好短期、有摊销机制的结构以对冲建设、租户和资产三重不确定性。不摊销的QTS债券利差比有施工保证的Hut 8宽约40bp，显示市场对资产到期价值的关注度上升。
{\small\color{refgray}数据：全球AI相关债务发行YTD 3360亿美元\par}
{\small\color{refgray}数据：高质量超大规模企业利差YTD走阔24bp\par}
{\small\color{refgray}数据：QTS债券利差比Hut 8宽约40bp\par}
{\small\color{refgray}边际变化：市场从全面投入转向有选择地关注，资产不确定性（续租、技术过时）成为定价核心变量，不摊销结构风险溢价上升。\par}
\paragraph{MS} 日本四大IT服务商（富士通、NEC、NRI、NTT数据）正用AI改写成本结构。外包成本占营收比从1991年的20\%升至2024年的33.5\%，成为最大支出。AI驱动开发使头部公司有能力将外包工作收回内部，富士通实现100倍效率提升（3人月→4小时），NEC实现99\%工作量削减。营业利润率有望在当前10-20\%基础上再提升5-10个百分点。
{\small\color{refgray}数据：外包成本占营收比33.5\%，超过人力成本27.3\%\par}
{\small\color{refgray}数据：富士通：3人月（480小时）压缩至4小时\par}
{\small\color{refgray}数据：NEC：6500小时→53小时（99\%削减）\par}
{\small\color{refgray}数据：NTT数据：2030年目标减少70\%开发流程\par}
{\small\color{refgray}边际变化：AI从概念验证进入规模化部署，外包成本结构性压缩成为利润改善最直接的杠杆，工程师角色从编码者转向AI监督者。\par}
\paragraph{UBS} 中国AI模型商业化分水岭到来，胜负变量从算力规模转向token ROI。编码能力成为模型智能自增强回路：代码可验证、可规模化生成训练数据，编码agent提供真实反馈。智谱ZCode、MiniMax Code、DeepSeek Code Harness正抢占飞轮入口。中国模型凭借结构性成本优势，在全球推理市场获得真实订单。
{\small\color{refgray}数据：编码agent已用于优化小模型训练代码和下一代模型迭代\par}
{\small\color{refgray}数据：智谱、MiniMax于2026年初上市\par}
{\small\color{refgray}数据：中国模型token成本低于美国竞品\par}
{\small\color{refgray}边际变化：模型层货币化TAM从开发者扩展到知识工作者，编码能力成为递归自我改进的早期迹象，SOTA模型领先优势可能比预期更持久。\par}
\paragraph{GS} 三菱重工与英伟达合作开发AI数据中心冷却与供电系统，短期业绩影响有限（数据中心业务占比<1\%），但增长路径清晰：2023年收购Concentric，2025年5月设立达拉斯基地，目标早期将业务规模提升至数千亿日元。市场尚未充分反映其在AI基础设施中的潜力。
{\small\color{refgray}数据：数据中心业务目前占三菱重工营收<1\%\par}
{\small\color{refgray}数据：2026财年增长领域（含数据中心）合计营收目标约1000亿日元\par}
{\small\color{refgray}数据：三菱重工目标早期数据中心业务达数千亿日元\par}
{\small\color{refgray}边际变化：AI研究焦点从GPU向基础设施（冷却、供电）扩展，三菱重工通过并购和基地建设为规模化铺路，估值尚未反映这一逻辑。\par}
\paragraph{GS} 欧洲电气化与AI数据中心推高电价的担忧被放大。未来十年家庭电费年均涨幅仅2-4\%，低于过去十年5\%的年增速。35-40\%的电气化投资（电网维护、陆上风电、光伏）非通胀性或具通缩效应；电力需求增长（电动车、热泵、数据中心）反而摊薄单位固定成本。德国和英国电费涨幅或达中高个位数，但整体可控。
{\small\color{refgray}数据：未来十年欧洲电气化投资2.2-3.5万亿欧元\par}
{\small\color{refgray}数据：35-40\%投资非通胀性或通缩性\par}
{\small\color{refgray}数据：2030年前电力需求年增速1.5-2.5\%，此后3-3.5\%\par}
{\small\color{refgray}数据：家庭全面电气化每年可节省约1200欧元燃料费用\par}
{\small\color{refgray}边际变化：市场线性推演大规模投资必然推高电价，但GS指出投资结构和需求增长两个因素可自然对冲电价上涨压力。\par}
\subsubsection*{关键数据}
\begin{itemize}
  \item 全球AI相关债务发行YTD 3360亿美元
  \item 高质量超大规模企业利差YTD走阔24bp
  \item 日本IT服务外包成本占营收33.5\%
  \item 富士通AI开发效率提升100倍（3人月→4小时）
  \item 中国AI模型token成本低于美国竞品
  \item 三菱重工数据中心业务占营收<1\%，目标数千亿日元
  \item 未来十年欧洲家庭电费年均涨幅2-4\%
  \item 欧洲电气化投资中35-40\%非通胀性
\end{itemize}
\begin{figure}[htbp]
\centering
\includegraphics[width=0.88\linewidth]{figures/fig_036.jpg}
\caption{Exhibit 1 - Exhibit 1: We're now at \textbackslash{}\$336bn of AI-related debt issuance YTD across global credit markets The terms below are largely the primary risk vectors associated with the data center construction debt that came up frequency in our inv}
\end{figure}
\begin{figure}[htbp]
\centering
\includegraphics[width=0.88\linewidth]{figures/fig_041.jpg}
\caption{Exhibit 1 - Exhibit 1: Outsourcing Cost Ratio to Revenue Japan Industry View In-Line MS does and seeks to do business with companies covered in MS. As a result, investors should be aware that the firm may have a conflict of interest that c}
\end{figure}
\begin{figure}[htbp]
\centering
\includegraphics[width=0.88\linewidth]{figures/fig_067.jpg}
\caption{Figure 1 - Figure 2: Chinese AI labs are increasingly launching first-party coding-agent products / harnesses}
\end{figure}
\begin{figure}[htbp]
\centering
\includegraphics[width=0.88\linewidth]{figures/fig_121.jpg}
\caption{Exhibit 1 - Exhibit 2: Germany and the UK may face the highest growth in electricity bills owing to rising grid fees}
\end{figure}
{\small\color{refgray}本节综合 5 篇研究；涉及 MS、UBS、GS。}
\newpage
\subsection{中国互联网与消费：利润优先转向，分化加剧}
\textbf{中国互联网与消费板块正经历从规模扩张到盈利优先的结构性转变，但各子行业和公司间的利润修复节奏、增长驱动力和竞争格局出现显著分化。}
\subsubsection*{主流共识}
\begin{itemize}
  \item 中国互联网公司普遍从追求规模转向盈利优先，利润修复成为核心叙事。
  \item 消费领域，龙头企业在弱市中通过产品结构优化和渠道调整巩固份额，利润韧性优于行业平均。
  \item AI云基础设施成为互联网公司增长新引擎，但应用层变现仍待突破。
\end{itemize}
\subsubsection*{分歧与条件差异}
\begin{itemize}
  \item 抖音与美团在本地生活的竞争：抖音核销率（57-58\%）远低于美团（>80\%），净交易额差距比毛交易额更大，但抖音正转向利润导向，目标2026下半年盈亏平衡。
  \item 京东二季度收入承压（零售收入同比-5\%），但利润端因品类结构改善和新业务减亏而超预期（non-GAAP净利润同比+13\%），市场对利润韧性的判断存在分歧。
  \item 五粮液去库存节奏慢于预期：终端动销双位数增长，但公司出货恢复滞后，2Q26净利润仅为2Q24的13-23\%，市场对库存周期结束时间看法不一。
  \item 化妆品行业利润分化显著：毛戈平、森林家园等领跑，华熙生物持续调整，线上ROI普遍走低，线下效率成为缓冲。
\end{itemize}
\subsubsection*{机构观点}
\paragraph{Bernstein} 市场对沃尔玛降价的担忧过度。沃尔玛的定价能力并非新闻，全渠道利润扩张才是核心：电商亏损收窄和零售媒体收入增长预计在FY26-FY30将美国业务EBIT利润率推高约200个基点至7\%。关税退税（约30亿美元）为降价提供缓冲，且沃尔玛杂货同店通胀率持续低于CPI食品通胀（价差从FY25 Q4的0\%扩大至FY27 Q1的1.9\%），已建立明确价差优势。
{\small\color{refgray}数据：沃尔玛美国EBIT利润率目标：FY30升至7\%（当前约5\%），电商亏损收窄贡献约100bps，零售媒体贡献约75bps\par}
{\small\color{refgray}数据：关税退税预计接近美国净销售额的50个基点，约30亿美元\par}
{\small\color{refgray}数据：沃尔玛美国杂货同店通胀率与CPI食品通胀价差：FY25 Q4为0\%，FY27 Q1扩大至1.9\%\par}
{\small\color{refgray}数据：零售媒体收入：当前约40亿美元（占GMV 4\%），预计GMV从1150亿增至2250亿时可达110亿美元\par}
{\small\color{refgray}边际变化：市场关注点从短期降价压力转向全渠道利润结构改善，关税退税提供了可预期的定价缓冲资金来源。\par}
\paragraph{MS} 京东2Q26业绩将好于市场担忧。预计京东零售收入同比-5\%（带电品类低双位数下滑），但经营利润率同比改善5bps，主要来自一般商品品类毛利率提升。集团non-GAAP净利润预计84亿元（同比+13\%），新业务亏损环比收窄（外卖亏损从70亿降至65亿，但被京喜和Joybuy投入抵消）。下半年收入增速因低基数触底回升，利润改善幅度更大。
{\small\color{refgray}数据：京东零售2Q26收入同比-5\%，其中电子与家电收入低双位数下滑\par}
{\small\color{refgray}数据：京东零售2Q26经营利润率同比改善5bps\par}
{\small\color{refgray}数据：集团2Q26 non-GAAP净利润预计84亿元，同比+13\%\par}
{\small\color{refgray}数据：新业务亏损：1Q26为104亿元，2Q26预计降至98亿元（外卖亏损从70亿降至65亿）\par}
{\small\color{refgray}边际变化：市场对京东二季度的担忧集中在收入端高基数，但利润端韧性（毛利率改善+新业务减亏）成为超预期来源。\par}
\paragraph{NOM} 中国互联网公司集体转向盈利优先，但抖音与美团在本地生活的竞争格局差异显著。抖音本地生活2026上半年毛GMV 5510亿元（同比+44\%），但核销率仅57-58\%，远低于美团的80\%以上，净交易额差距更大。抖音正调整策略：酒店旅行成为增长主力，加强佣金收取，目标2026下半年盈亏平衡。AI大模型进入闭源变现阶段，月之暗面年化经常性收入有望年底前突破10亿美元。
{\small\color{refgray}数据：抖音本地生活2026上半年毛GMV 5510亿元，同比+44\%\par}
{\small\color{refgray}数据：抖音本地生活核销率57-58\%，美团到店核销率>80\%\par}
{\small\color{refgray}数据：月之暗面年化经常性收入目标：年底前突破10亿美元\par}
{\small\color{refgray}数据：抖音本地生活目标：2026下半年实现盈亏平衡\par}
{\small\color{refgray}边际变化：抖音从规模扩张转向利润导向，核销率差异成为理解与美团竞争的关键变量；AI大模型从开源转向闭源变现。\par}
\paragraph{GS} 百度2Q26收入结构发生实质性转变：AI云业务首次超过传统广告成为百度核心收入最大来源。预计AI云收入同比+41\%（GPU云三位数增长），广告收入同比-21\%。核心经营利润同比-18\%，与1Q26的-19\%基本企稳，得益于AI云利润率逐季改善。但为追赶AI模型训练，下半年资本开支将加速，持续影响利润修复。
{\small\color{refgray}数据：百度核心2Q26收入预计252亿元，同比-4\%\par}
{\small\color{refgray}数据：AI云收入同比+41\%，其中AI云基础设施（GPU云）同比+58\%\par}
{\small\color{refgray}数据：广告收入同比-21\%，占核心收入比例从2025年>60\%降至2026年<50\%\par}
{\small\color{refgray}数据：核心经营利润同比-18\%（1Q26为-19\%），呈企稳态势\par}
{\small\color{refgray}边际变化：百度增长引擎从广告切换至AI云基础设施，但资本开支加速将压制利润修复节奏。\par}
\paragraph{GS} 中国化妆品公司2Q26收入增长加速，但利润分化加剧。毛戈平（1H26收入+28\%）和森林家园（+39\%）领跑，珀莱雅增速温和（+1\%），巨子生物受高基数影响（-7\%），华熙生物继续调整（-18\%）。新兴品牌（Off-relax、花知晓等）成为增长主力，最高可实现50\%同比增长。线上利润率整体呈节奏变化，线下效率成为缓冲。
{\small\color{refgray}数据：毛戈平1H26收入同比+28\%，森林家园+39\%\par}
{\small\color{refgray}数据：上海家化2Q26收入+14\%，贝泰妮+10\%，珀莱雅+1\%\par}
{\small\color{refgray}数据：巨子生物1H26收入同比-7\%，华熙生物-18\%\par}
{\small\color{refgray}数据：新兴品牌最高可实现50\%同比增长\par}
{\small\color{refgray}边际变化：增长重心从单一主力品牌转向多品牌矩阵，线上ROI走低背景下线下效率成为利润缓冲。\par}
\paragraph{GS} 九号公司2Q26收入增长加速（同比+27\%，1Q26为+15\%），利润不确定性环比大幅收窄（净利润同比-5\%，1Q26为-55\%）。割草机器人弹性较高增长（毛利率>48\%），收入占比从2024年6.3\%提升至2026年预期12.6\%，推动整体利润率修复。两轮车需求走出低谷（6月行业降幅收窄至-4\%），公司持续获取份额。汇率和成本压力预计Q2见顶，下半年逐季缓解。
{\small\color{refgray}数据：九号公司2Q26收入同比+27\%（1Q26为+15\%）\par}
{\small\color{refgray}数据：2Q26净利润同比-5\%（1Q26为-55\%），剔除汇兑损益后经常性利润双位数增长\par}
{\small\color{refgray}数据：割草机器人业务毛利率>48\%，收入占比从2024年6.3\%升至2026年预期12.6\%\par}
{\small\color{refgray}数据：国内两轮车行业6月降幅收窄至-4\%（1Q26为-17\%）\par}
{\small\color{refgray}边际变化：成本上涨和汇率波动冲击见顶，高毛利割草机器人占比提升成为利润修复核心驱动力。\par}
\paragraph{GS} 五粮液1H26净利润预告87.3-92亿元（同比+89-99\%），但仅为1H24的46-48\%。2Q26隐含净利润中值约9亿元，仅为2Q24的18\%，远低于GS预期的25亿元。终端动销双位数增长，但公司出货恢复滞后，渠道去库存仍在进行。低基数效应减弱，下半年同比增速可能进一步收窄。
{\small\color{refgray}数据：五粮液1H26净利润中值约89.7亿元，同比+94\%，但较GS预期低15\%\par}
{\small\color{refgray}数据：2Q26净利润中值约9亿元，仅为2Q24（50亿元）的18\%\par}
{\small\color{refgray}数据：普五终端动销年初至今双位数增长，但公司出货恢复明显滞后\par}
{\small\color{refgray}数据：对应16倍2027E PE，目标价70元\par}
{\small\color{refgray}边际变化：利润反弹弱于预期，渠道去库存周期比市场预期更持久，终端动销与公司出货的差距是核心观察指标。\par}
\paragraph{JPM} 中国家电六月零售数据（同比-19\%）不宜过度解读，主要受补贴退坡、天气偏凉和618促销节奏变化等噪音因素影响。两年复合增速下白电销量基本持平（空调零增长），需求走平而非走弱。均价端承压（白电均价下降中高个位数），消费者倾向经济型产品。份额格局在弱市中巩固：美的空调份额34.9\%（+1.7ppt），海尔冰箱份额43.7\%（+3.7ppt）。
{\small\color{refgray}数据：六月家电零售额同比-19\%（五月为-9\%），两年复合增速转负至-4\%\par}
{\small\color{refgray}数据：空调六月销售额同比-28\%，占六月销售超40\%\par}
{\small\color{refgray}数据：两年复合增速下白电销量基本持平，空调零增长\par}
{\small\color{refgray}数据：白电均价两年复合下降中高个位数\par}
{\small\color{refgray}边际变化：单月数据噪音较大，两年复合视角下需求走平而非走弱，龙头份额在弱市中进一步巩固。\par}
\subsubsection*{关键数据}
\begin{itemize}
  \item 沃尔玛美国EBIT利润率目标FY30升至7\%，电商亏损收窄+零售媒体贡献约175bps
  \item 京东2Q26 non-GAAP净利润84亿元，同比+13\%，零售OPM同比改善5bps
  \item 抖音本地生活核销率57-58\% vs 美团>80\%，净GMV差距更大
  \item 百度AI云收入同比+41\%，首次超过广告成为核心收入最大来源
  \item 毛戈平1H26收入+28\%，森林家园+39\%，华熙生物-18\%
  \item 九号公司2Q26收入+27\%，净利润-5\%（1Q26为-55\%），割草机器人毛利率>48\%
  \item 五粮液2Q26净利润仅为2Q24的18\%，去库存慢于预期
  \item 六月家电零售同比-19\%，两年复合增速-4\%，白电均价下降中高个位数
\end{itemize}
\begin{figure}[htbp]
\centering
\includegraphics[width=0.88\linewidth]{figures/fig_087.jpg}
\caption{Exhibit 3 - Exhibit 3: Baidu: Core AI-powered business quarterly revenue and yoy \% Exhibit 4: Baidu: AI Cloud Infra quarterly revenue and yoy \% \#\# Financials}
\end{figure}
\begin{figure}[htbp]
\centering
\includegraphics[width=0.88\linewidth]{figures/fig_149.jpg}
\caption{Figure 1 - Figure 1: China home appliance - AVC retail sales YoY\% Average of online and offline retail sales YoY\% Figure 2: Air con - retail sales YoY\% Average of online and offline retail sales YoY\% Figure 3: Refrigerator - retail sales}
\end{figure}
\begin{figure}[htbp]
\centering
\includegraphics[width=0.88\linewidth]{figures/fig_020.jpg}
\caption{EXHIBIT 3 - EXHIBIT 3: As a result, Walmart US has been taking share in recent years... Mass Retailer YoY Market Share Gain/Loss vs. Core Retail Sales (bps) EXHIBIT 4: ...with more pronounced share gains at Sam's Club. Club Store YoY Market}
\end{figure}
\begin{figure}[htbp]
\centering
\includegraphics[width=0.88\linewidth]{figures/fig_089.jpg}
\caption{Exhibit 7 - Exhibit 7: We use 22X as MA base multiple, which is based on trough industry P/E between Jan 25 - Jul 26 China MA Names average P/E NTM PE: China Cosmetics Coverage Companies include Proya, Botanee, Shanghai Jahwa, Giant Biogene,}
\end{figure}
{\small\color{refgray}本节综合 8 篇研究；涉及 Bernstein、MS、NOM、GS、JPM。}
\newpage
\subsection{汽车与新能源：全球分化加剧，中国品牌海外扩张成核心主线}
\textbf{全球汽车与新能源市场呈现显著区域分化：中国纯电渗透率创新高但国内需求疲弱，欧洲受政策与平价车型驱动增长强劲，美国电动化进程遇阻。中国品牌凭借海外市场超预期扩张成为结构性亮点，盈利修正方向成为选股核心逻辑。}
\subsubsection*{主流共识}
\begin{itemize}
  \item 中国新能源渗透率持续提升，6月纯电渗透率达43\%，整体新能源渗透率63\%，电动化趋势未逆转。
  \item 全球纯电销量6月环比增长11\%至104.5万辆，累计同比缺口收窄至仅1\%，市场正在走出年初低谷。
  \item 中国品牌在欧洲市场份额持续扩大，5月已达13\%，较2025年的9\%显著提升，海外销量指引存在20-50\%上修空间。
\end{itemize}
\subsubsection*{分歧与条件差异}
\begin{itemize}
  \item 国内需求前景：JPM认为下半年零售仍将低于季节性趋势，全年跌幅约15\%；而GS周度数据显示7月首周零售同比降幅已从-20\%收窄至-15\%，边际改善。
  \item 美国电动化路径：JPM数据显示6月纯电渗透率降至5\%，累计同比-26\%；但GS认为现代起亚凭借混动产品在美国实现净价格同比+2.6\%，份额持续扩大，混动路线可能优于纯电。
  \item 估值逻辑切换：GS认为现代摩比斯估值扩张来自人形机器人执行器业务（2035年贡献7\%利润），而非传统零部件；JPM则强调中国车企盈利修正方向比估值更重要。
\end{itemize}
\subsubsection*{机构观点}
\paragraph{GS} 中国新能源车市场进入产品力驱动分化阶段。7月第二周整体订单环比下降16\%，但蔚来（+13\%）、零跑（+10\%）、小鹏（+5\%）实现正增长，蔚来五座版ES8是短期催化剂。乘用车零售同比降幅收窄至-15\%，新能源渗透率60.5\%。电池级碳酸锂价格周跌7.3\%至15.3万元/吨，电芯价格稳定。
{\small\color{refgray}数据：7月第二周整体新能源订单环比-16\%，同比-19\%\par}
{\small\color{refgray}数据：蔚来周订单环比+13\%，同比+114\%\par}
{\small\color{refgray}数据：蔚来2026年累计订单同比+98\%\par}
{\small\color{refgray}数据：7月1-5日乘用车零售169k辆，同比-15\%\par}
{\small\color{refgray}边际变化：从6月底新品集中上市高峰自然回落，但蔚来等品牌逆势增长，市场从'新品脉冲'转向'产品力驱动'分化阶段。\par}
\paragraph{GS} 现代摩比斯正从汽车零部件供应商转型为人形机器人执行器制造商，估值逻辑切换。公司间接持有波士顿动力11.25\%股权，但估值扩张更依赖执行器业务从零到一的利润贡献：2027年起贡献收入，2035年预计贡献公司7\%营业利润，营业利润率18\%。当前约9倍远期PE仍低于全球零部件同业平均11倍。
{\small\color{refgray}数据：摩比斯间接持有波士顿动力11.25\%股权\par}
{\small\color{refgray}数据：当前约9倍一年期远期PE，低于全球同业11倍\par}
{\small\color{refgray}数据：2027年执行器收入预计6340亿韩元，占比1\%\par}
{\small\color{refgray}数据：2035年执行器收入预计2.6万亿韩元，占比3\%\par}
{\small\color{refgray}边际变化：市场从仅关注波士顿动力股权重估，转向将摩比斯视为执行器量产平台，估值叙事从'汽车零部件'切换至'机器人执行器'。\par}
\paragraph{GS} 现代起亚在美国市场展现定价韧性。6月组合净价格34,717美元，同比+2.6\%，显著高于行业平均同比+0.5\%。混动产品线是定价能力的结构性来源。6月现代美国市占率6.2\%，起亚5.1\%，GS预计2026年分别升至6.4\%和5.6\%，较2025年各提升0.4个百分点。
{\small\color{refgray}数据：现代起亚6月组合净价格34,717美元，同比+2.6\%\par}
{\small\color{refgray}数据：行业6月净价格46,294美元，同比+0.5\%\par}
{\small\color{refgray}数据：现代起亚净价格同比增速高出行业2.1个百分点\par}
{\small\color{refgray}数据：6月现代美国市占率6.2\%，起亚5.1\%\par}
{\small\color{refgray}边际变化：在行业价格承压背景下，现代起亚凭借混动产品组合维持净价格同比正增长，定价能力差异扩大，份额扩张获得数据支撑。\par}
\paragraph{JPM} 中国汽车行业下半年选股核心逻辑是盈利修正方向。1H26 MSCI中国汽车指数下跌17\%，但中国重汽（+49\%）、吉利（+4\%）、蔚来（-6\%）因盈利上修表现相对领先。国内需求短期难复苏，1H26零售同比-23\%，下半年预计仍低于季节性。海外市场是结构性增长来源，中国品牌欧洲市占率已达13\%，海外销量指引存在20-50\%上修空间。
{\small\color{refgray}数据：1H26 MSCI中国汽车指数下跌17\%，MSCI中国指数下跌10\%\par}
{\small\color{refgray}数据：中国重汽1H26上涨49\%，吉利上涨4\%，蔚来下跌6\%\par}
{\small\color{refgray}数据：1H26中国乘用车零售同比下滑23\%\par}
{\small\color{refgray}数据：JPM预计2026全年零售跌幅约15\%\par}
{\small\color{refgray}边际变化：从关注销量和估值转向盈利修正方向，海外市场成为区分盈利质量的核心变量，欧洲市占率超预期提升。\par}
\paragraph{JPM} 全球纯电销量6月环比增长11\%至104.5万辆，同比+9\%，累计同比缺口收窄至仅1\%。区域分化显著：中国纯电渗透率升至43\%创新高，EU-10环比+34\%且同比+48\%为最强增长引擎，美国纯电渗透率降至5\%且累计同比-26\%表现最弱。锂辉石价格短期偏软，但中期受储能需求支撑。
{\small\color{refgray}数据：6月全球纯电销量104.5万辆，环比+11\%，同比+9\%\par}
{\small\color{refgray}数据：中国6月纯电销量68.5万辆，环比+8\%，渗透率43\%\par}
{\small\color{refgray}数据：中国6月整体新能源渗透率63\%，与2025年8月高点持平\par}
{\small\color{refgray}数据：EU-10 6月纯电销量28.7万辆，环比+34\%，同比+48\%\par}
{\small\color{refgray}边际变化：全球纯电销量从年初低谷修复，但区域分化加剧：欧洲成为最强增长引擎（政策+平价车型双轮驱动），美国电动化进程遇阻（渗透率不升反降）。\par}
\subsubsection*{关键数据}
\begin{itemize}
  \item 7月第二周中国新能源整体订单环比-16\%，蔚来环比+13\%
  \item 蔚来2026年累计订单同比+98\%
  \item 现代起亚美国6月净价格同比+2.6\%，行业仅+0.5\%
  \item 现代摩比斯当前约9倍远期PE，低于全球同业11倍
  \item 中国品牌欧洲市占率5月达13\%（2025年为9\%）
  \item 全球6月纯电销量104.5万辆，环比+11\%，同比+9\%
  \item 中国6月纯电渗透率43\%，美国降至5\%
  \item EU-10 6月纯电销量环比+34\%，同比+48\%
\end{itemize}
\begin{figure}[htbp]
\centering
\includegraphics[width=0.88\linewidth]{figures/fig_136.jpg}
\caption{Exhibit 1 - Exhibit 1: HMK's US market share vs. net price growth spread over the industry Exhibit 2: US net price (ATP - incentives) trend Exhibit 3: HMC \& Kia's YoY net price growth vs. industry}
\end{figure}
\begin{figure}[htbp]
\centering
\includegraphics[width=0.88\linewidth]{figures/fig_144.jpg}
\caption{Figure 1 - Figure 1: China autos - Domestic (retail) PV demand growth analysis and projection in 2026 vs. historical seasonality (we project domestic demand may largely fall below seasonality in 2H26)}
\end{figure}
\begin{figure}[htbp]
\centering
\includegraphics[width=0.88\linewidth]{figures/fig_155.jpg}
\caption{Figure 1 - Figure 1: China vehicle sales vs. BEV sales Figure 2: China EV penetration rates Battery Electric Vehicles (\%) Plug-in Hybrid Electric Vehicles (\%) Total EV (BEV+PHEV) penetration (\%) \#\# Head of Australia Metals \& Mining}
\end{figure}
\begin{figure}[htbp]
\centering
\includegraphics[width=0.88\linewidth]{figures/fig_157.jpg}
\caption{Figure 3 - Figure 3: Battery electric vehicle penetration rates by region Figure 4: Battery electric vehicle sales – monthly (China + EU + US) '000 units Figure 5: Global vehicle sales}
\end{figure}
{\small\color{refgray}本节综合 5 篇研究；涉及 GS、JPM。}
\newpage
\subsection{医疗健康：GLP-1主导代谢赛道，心血管与老龄化开启新叙事}
\textbf{GLP-1市场2031年将达千亿美元，礼来占据约70\%份额，口服药并非颠覆者；心血管领域Lp(a)靶点面临关键验证，老龄化推动后急性期护理需求结构性增长；生命科学工具板块温和回暖，中国ICL行业供给出清带来利润弹性。}
\subsubsection*{主流共识}
\begin{itemize}
  \item GLP-1市场2031年规模约1020亿美元，礼来凭借Mounjaro/Zepbound占据约70\%份额，注射剂长期主导（2030年后>79\%）。
  \item 生命科学工具终端市场连续两个季度环比改善，工业端增速达高个位数，接近历史高位但仍有上行空间。
  \item 婴儿潮一代进入80岁将推动医疗支出从预防性向长期照护迁移，85岁以上人均支出约36000美元，是65-74岁的1.6倍。
  \item Lp(a)是最后一个未成药的主要心血管风险因子，HORIZON试验若成功将验证靶点并重塑赛道估值。
\end{itemize}
\subsubsection*{分歧与条件差异}
\begin{itemize}
  \item 口服GLP-1份额：Bernstein认为2029年达22\%峰值后回落，但市场对口服药颠覆性预期更高，且Foundayo在T2D领域被低估。
  \item Lp(a)降低幅度与获益关系：GS认为线性关系下安进olpasiran（降幅\textasciitilde{}105\%）优于诺华pelacarsen（\textasciitilde{}50\%），但临床最优降幅仍存争议。
  \item 中国ICL行业：GS认为金域医学利润修复靠成本管控而非收入反弹，但市场担忧DRG 3.0量价齐跌。
  \item 心血管新药催化剂：GS对HORIZON/ZEUS的波动范围测算（NVS +2\%/-3\%，NVO +7\%/-5\%）低于部分投资者预期。
\end{itemize}
\subsubsection*{机构观点}
\paragraph{Bernstein} 美国GLP-1品牌市场2031年规模约1020亿美元，礼来占据约70\%份额。口服药份额2029年达22\%峰值后回落，注射剂长期主导（>79\%）。Foundayo约47\%的2031年美国收入来自T2D，该领域被市场低估。患者渗透率持续加深：2031年约50\% T2D患者使用肠促胰素（目前25\%），肥胖/超重渗透率升至23\%。
{\small\color{refgray}数据：2031年GLP-1市场1020亿美元\par}
{\small\color{refgray}数据：礼来占约70\%份额\par}
{\small\color{refgray}数据：口服药份额2029年峰值22\%\par}
{\small\color{refgray}数据：Foundayo 47\%收入来自T2D\par}
{\small\color{refgray}边际变化：下调Foundayo美国预期，因更温和上市节奏；口服药峰值份额从更高预期修正至22\%。\par}
\paragraph{UBS} 2026年首批婴儿潮一代满80岁，是医疗服务需求结构根本性转变的起点。80岁以上人口将从2020年代中期的约1500万增至2040年代的近3000万。85岁以上人均医疗支出约36000美元，比65-74岁高60\%以上。护理模式从治疗转向时长，后急性期护理（专业护理、居家护理、临终关怀）成为主战场。
{\small\color{refgray}数据：2026年首批婴儿潮满80岁\par}
{\small\color{refgray}数据：80+人口从1500万增至3000万（2040s）\par}
{\small\color{refgray}数据：85+人口2040年达1400万（+72\% vs 2022）\par}
{\small\color{refgray}数据：85+人均支出36000美元\par}
{\small\color{refgray}边际变化：从泛老龄化叙事聚焦至80岁分水岭，强调80岁后医疗支出从预防性向长期照护的结构性迁移。\par}
\paragraph{GS} 生命科学工具终端市场连续两个季度环比改善，GS Tools Tracker显示2Q26增速较1Q26轻微加速。工业端增速达高个位数（长期均值低个位数），PMI和半导体资本支出环比提升是主要驱动，按历史规律仍有上行空间。生物制药端低两位数增长但略有放缓，临床前VC融资传导至资本开支需1年以上。中国设备更新订单2Q26同比降57\%，但布鲁克等受影响有限。
{\small\color{refgray}数据：2Q26工业端增速高个位数\par}
{\small\color{refgray}数据：工业端长期均值低个位数\par}
{\small\color{refgray}数据：中国设备更新订单同比降57\%\par}
{\small\color{refgray}数据：生物制药端低两位数增长\par}
{\small\color{refgray}边际变化：从市场担忧结构性放缓转向确认后疫情余波，连续两个季度加速信号改善下半年有机增长预期。\par}
\paragraph{GS} CARDIO-TTRansform试验失败后，Q3两个关键催化剂可能重塑心血管赛道：NVS/IONS的HORIZON试验（Lp(a)靶点首次大规模验证，n=8323）和NVO的ZEUS研究（IL-6抗体降低炎症）。HORIZON整体人群需12\%风险降低达统计显著，高风险亚组需14\%。若成功，将验证Lp(a)靶点并提升安进olpasiran、礼来muvalaplin等资产估值。
{\small\color{refgray}数据：Lp(a)影响全球约14亿人\par}
{\small\color{refgray}数据：HORIZON入组8323人\par}
{\small\color{refgray}数据：整体人群需12\%风险降低\par}
{\small\color{refgray}数据：高风险亚组需14\%风险降低\par}
{\small\color{refgray}边际变化：从CARDIO-TTRansform失败后的悲观情绪转向关注HORIZON/ZEUS两个潜在催化剂，量化了不同结果对相关公司的波动范围。\par}
\paragraph{GS} 金域医学1H26净利润1.7-2.2亿元（去年同期净亏损8500万元），超预期主要来自成本控制与数字化转型。股权激励目标隐含管理层对DRG 3.0和价格标准化的预期：2026-2028年营收年增6-10\%，扣非净利润从3.9亿增至6亿，利润增速快于收入。行业供给出清后利润率弹性大于收入弹性。
{\small\color{refgray}数据：1H26净利润1.7-2.2亿元\par}
{\small\color{refgray}数据：去年同期净亏损8500万元\par}
{\small\color{refgray}数据：2026/27/28年营收增速≥6\%/14.5\%/26\%\par}
{\small\color{refgray}数据：扣非净利润目标3.9/5/6亿元\par}
{\small\color{refgray}边际变化：从后疫情利润压缩转向确认内部效率改善驱动的利润修复，股权激励目标显示管理层对政策不确定性有预期且利润弹性更大。\par}
\subsubsection*{关键数据}
\begin{itemize}
  \item 2031年GLP-1市场1020亿美元，礼来占70\%
  \item 口服药份额2029年峰值22\%，注射剂长期>79\%
  \item 85+人均医疗支出36000美元，比65-74岁高60\%
  \item 80+人口从1500万增至3000万（2040s）
  \item 生命科学工具2Q26工业端增速高个位数
  \item 中国设备更新订单2Q26同比降57\%
  \item HORIZON试验n=8323，需12\%风险降低达显著
  \item 金域医学1H26净利润1.7-2.2亿元（扭亏）
  \item 股权激励目标：扣非净利润从3.9亿增至6亿（2026-2028）
\end{itemize}
\begin{figure}[htbp]
\centering
\includegraphics[width=0.88\linewidth]{figures/fig_009.jpg}
\caption{EXHIBIT 2 - EXHIBIT 2: Revenue (\textbackslash{}\$m), split by Company EXHIBIT 4: We see Mounjaro retaining majority share of Type 2 Diabetes (T2D), GLP-1 market out to 2035}
\end{figure}
\begin{figure}[htbp]
\centering
\includegraphics[width=0.88\linewidth]{figures/fig_077.jpg}
\caption{Figure 1 - Figure 2: Prevalence of Chronic Conditions Increases with Age}
\end{figure}
\begin{figure}[htbp]
\centering
\includegraphics[width=0.88\linewidth]{figures/fig_091.jpg}
\caption{Exhibit 1 - Exhibit 1: Continued bifurcation of end market growth according to the GS Tools Tracker in 2Q26 Exhibit 2: mAb Sales by volume tracking above average historical growth in recent months KG = Kilograms Exhibit 3: Preclinical Biot}
\end{figure}
\begin{figure}[htbp]
\centering
\includegraphics[width=0.88\linewidth]{figures/fig_131.jpg}
\caption{Exhibit 2 - Exhibit 2: Analysis of the placebo arm in ODYSSEY OUTCOMES highlights that LDL-C, hsCRP and Lp(a) are independent risk factors for MACE, but Lp(a) appeared to show limited correlation with all cause and CV death A Mechanistically}
\end{figure}
{\small\color{refgray}本节综合 6 篇研究；涉及 Bernstein、UBS、GS。}
\newpage
\subsection{工业与材料：供给出清与结构分化下的新定价锚}
\textbf{工业与材料板块正经历从总量驱动到结构驱动的切换：化工行业供给端出清信号增强，建材翻新需求接棒新建，造船订单向中国集中，而AI与能源转型则催生设备与材料的新增量。市场定价逻辑从宏观预期转向微观兑现，个股分化加剧。}
\subsubsection*{主流共识}
\begin{itemize}
  \item 化工行业供给端出清加速，亚洲乙烯价差已跌至触发减产的水平，中国反内卷政策与韩国裂解装置减产共同推动行业重组。
  \item 建材行业总量企稳，翻新需求与二手房交易回升正在对冲新建面积下滑，长期TAM有望回到2021年峰值。
  \item AI基础设施投资正从基建转向设备采购，PCB设备、电子化学品等细分领域受益明确。
  \item 中国造船业全球份额持续提升，6月按CGT计达85\%，在集装箱船、油轮、散货船三大主力船型上占据主导地位。
\end{itemize}
\subsubsection*{分歧与条件差异}
\begin{itemize}
  \item 数据中心供电路径分歧：BTM模式（Bloom）交付周期55-90天但LCOE高达\$82-115/MWh，FTM模式（Fervo/Ormat）成本更低但需等待2-4年电网接入，市场对速度与规模的定价存在差异。
  \item 化工板块反弹持续性存疑：氮肥与甲醇等供给驱动型品种年初至今涨幅超50\%，但工业化学品仍受宏观周期压制，盈利兑现能力有待验证。
  \item eVTOL商业化节奏分歧：政策支持力度大但适航审批周期超预期，GS将亿航2027年营收预测下调12\%，市场对时间假设的校准尚未完成。
  \item 中国船厂内部竞争格局分化：恒力重工凭借MSC大单加速扩张，扬子江船业接单节奏低于Q1均值24\%，产能扩张与产品结构升级带来差异化表现。
\end{itemize}
\subsubsection*{机构观点}
\paragraph{Bernstein} 数据中心供电存在结构性供给缺口，传统电网扩容需3-8年，而超大规模客户需求按季度计算。Bloom Energy的BTM模式（燃料电池直接部署在客户园区）交付周期仅55-90天，完全绕过电网互联瓶颈，但LCOE高达\$82-115/MWh，包含天然气燃料成本和每5-7年约\$1000/kW的电池堆更换储备。Fervo Energy的FTM模式（地热发电接入电网）建设成本约\$7000/kW，目标降至\$3000/kW，但需2-4年通电。Ormat Technologies同样采用FTM模式。短期（2026-2028）BTM模式在速度上占优，但长期成本曲线决定规模。
{\small\color{refgray}数据：Bloom Energy交付周期55-90天\par}
{\small\color{refgray}数据：Bloom LCOE \$82-115/MWh\par}
{\small\color{refgray}数据：Fervo建设成本\$7000/kW，目标\$3000/kW\par}
{\small\color{refgray}数据：传统电网互联需3-4年\par}
{\small\color{refgray}边际变化：首次将三家公司的数据中心供电合同纳入统一分析框架（BTM/FTM/BYOG），揭示不同商业模式在速度、成本和不确定性暴露上的根本差异。\par}
\paragraph{MS} 日本化工行业正站在供给端出清的拐点上。石化板块行业观点定为'Attractive'，核心证据：亚洲乙烯利用率持续承压但价差进一步下行空间有限；中国反内卷政策与韩国石脑油裂解装置减产信号叠加，行业重组加速；板块估值处于低位，情绪改善来自供给端出清预期而非需求复苏。电子化学品受益于AI半导体扩张和300mm硅片恢复趋势，增长路径更可预期。精细化学品依赖航空复材全面恢复。
{\small\color{refgray}数据：亚洲乙烯利用率持续承压\par}
{\small\color{refgray}数据：韩国出现石脑油裂解装置减产迹象\par}
{\small\color{refgray}数据：300mm硅片恢复趋势明显\par}
{\small\color{refgray}边际变化：从'需求复苏驱动'转向'供给端出清驱动'，强调中国和韩国的产能调整信号比市场预期更早到来。\par}
\paragraph{UBS} 大族激光2026年上半年净利润同比增长156\%-177\%，超出UBS和路透一致预期的109\%。增长并非依赖单一客户，而是PCB设备、消费电子设备、新能源、泛半导体和通用激光五条业务线同时发力。PCB设备收入同比增长超100\%，子公司大族数控净利润率从2025年Q2的10\%提升至2026年Q2的约20\%。AI PCB厂商资本开支正从基建转向高价值设备采购，胜宏科技和Avary 2026年资本开支指引分别同比增长172\%和154\%，行业新增AI PCB投资超1000亿元。公司正从设备供应商升级为AI基础设施核心工艺方案商。
{\small\color{refgray}数据：H126净利润同比+156\%至177\%\par}
{\small\color{refgray}数据：PCB设备收入同比+100\%\par}
{\small\color{refgray}数据：大族数控净利率从10\%升至20\%\par}
{\small\color{refgray}数据：胜宏科技2026年资本开支指引同比+172\%\par}
{\small\color{refgray}边际变化：市场认知从'苹果产业链周期股'升级为'AI基础设施核心工艺方案商'，超快激光钻孔设备在高端HDI和mSAP工艺中的应用重塑增长曲线。\par}
\paragraph{GS} 2021-2025年中国建材行业TAM调整约20\%，但四家覆盖公司市值跌幅达70\%，估值与基本面背离。行业总量正在企稳，驱动逻辑从新建转向翻新：积压的翻新需求与二手房交易回升对冲新建面积下滑，长期TAM有望回到2021年峰值。四家公司未来三年收入增长几乎全部来自零售渠道扩张、低线市场渗透和产品线拓宽。防水材料和涂料价格已调整，原材料成本回落，价差扩大，叠加零售占比提升，预计到2028年毛利率提升1-3个百分点。2026-2028年EPS复合增速9\%-29\%。
{\small\color{refgray}数据：2021-2025年行业TAM调整约20\%\par}
{\small\color{refgray}数据：四家公司市值跌幅达70\%\par}
{\small\color{refgray}数据：预计2028年毛利率提升1-3个百分点\par}
{\small\color{refgray}数据：2026-2028年EPS复合增速9\%-29\%\par}
{\small\color{refgray}边际变化：从'新建驱动'转向'翻新驱动'，渠道结构优化（从大B开发商转向C端/中小B端）改变定价权和回款周期，估值体系需要重新评估。\par}
\paragraph{GS} 中国低空经济政策密集落地，但亿航商业化节奏因监管审批周期超出预期而被重新审视。EH216-S商业载客运营仍需等待民航局新的运行安全要求。GS将2027年营收预测下调12\%，2028年下调8\%，主因AAV出货量假设过于激进。应用场景分三层：短期空中观光（点A到点A），中期短途物流和机场接驳（点A到点B），长期空中出租车（按需飞行）。物流场景推进可能快于载人。地方补贴力度大：安徽新航线奖励40万元，深圳每架次载人飞行补贴100-300元，广州单条载人航线运营补贴上限100万元。
{\small\color{refgray}数据：2027年营收预测下调12\%\par}
{\small\color{refgray}数据：2028年营收预测下调8\%\par}
{\small\color{refgray}数据：安徽新航线补贴40万元/条\par}
{\small\color{refgray}数据：深圳每架次载人飞行补贴100-300元\par}
{\small\color{refgray}边际变化：估值锚点从'政策利好'转向'审批节点可预测性'，时间假设修正导致目标价从\$16.9调整至\$7.3。\par}
\paragraph{GS} 6月全球新船订单量环比降9\%，同比微增1\%，但订单金额环比降39\%，主因高价值气体运输船订单减少（LPG/LNG船订单量环比降60-70\%）。三大主力船型（集装箱船、油轮、散货船）订单量价合计环比增长22\%和14\%。中国船厂全球订单份额（按CGT计）升至85\%，创2024年8月以来新高；剔除LNG船后升至87\%，韩国份额从40\%降至9\%。恒力重工6月新订单量较4-5月均值增长102\%，MSC向其订购最多20艘2万TEU双燃料集装箱船，估值约40亿美元。
{\small\color{refgray}数据：6月全球新船订单量环比-9\%，同比+1\%\par}
{\small\color{refgray}数据：订单金额环比-39\%\par}
{\small\color{refgray}数据：LPG/LNG船订单量环比-60\%至70\%\par}
{\small\color{refgray}数据：中国船厂份额85\%（按CGT计）\par}
{\small\color{refgray}边际变化：中国船厂份额从5月的约70\%跳升至85\%，恒力重工凭借MSC大单改写竞争格局，扬子江面临更激烈竞争。\par}
\paragraph{JPM} 2026年初至今北美化工板块内部分化显著：Compass Minerals、CF Industries、Chemours、Tronox等2025年调整较深的公司年初至今涨幅均超47\%，而Sherwin-Williams、Ecolab、Linde等防御性标的涨幅仅3\%-24\%。氮肥与甲醇成为反弹先锋，主因全球粮食价格维持高位叠加供给偏紧。包装板块温和回暖，Ball Corp上涨16.5\%、Silgan上涨11.4\%。工业化学品（LyondellBasell、Dow）虽年初至今涨超30\%，但受宏观周期影响波动较大。定价逻辑正从'超跌反弹'切换至'盈利兑现'阶段。
{\small\color{refgray}数据：Compass Minerals年初至今+51.9\%\par}
{\small\color{refgray}数据：CF Industries年初至今+51.2\%\par}
{\small\color{refgray}数据：标普500年初至今+10.7\%\par}
{\small\color{refgray}数据：FMC年初至今-21.3\%\par}
{\small\color{refgray}边际变化：市场从对宏观不确定性的集中定价转向重新审视各公司个体叙事，供给端故事（氮肥、甲醇）成为反弹核心驱动。\par}
\paragraph{JPM} 2026年7月9日发布的十五五碳达峰行动方案，核心指标为2030年单位GDP碳排放较2025年下降17\%、非化石能源占比达25\%。真正新信息在执行层面：每年非化石能源占比提升约1个百分点、煤电效率达标容量提升约15个百分点、零碳矿山/园区试点。这些量化抓手意味着电网设备、储能和新能源相关企业进入订单可见性更高的阶段。中广核电力二季度核电发电量586.40亿千瓦时（同比+3.5\%），长江电力二季度水电发电量709亿千瓦时（同比+2.8\%）。
{\small\color{refgray}数据：2030年非化石能源占比目标25\%\par}
{\small\color{refgray}数据：每年非化石能源占比提升约1个百分点\par}
{\small\color{refgray}数据：煤电效率达标容量提升约15个百分点\par}
{\small\color{refgray}数据：中广核电力Q2核电发电量586.40亿千瓦时，同比+3.5\%\par}
{\small\color{refgray}边际变化：政策从'定目标'转向'定节奏'，量化执行机制使电网设备、储能和新能源相关企业订单可见性显著提升。\par}
\subsubsection*{关键数据}
\begin{itemize}
  \item Bloom Energy交付周期55-90天，LCOE \$82-115/MWh
  \item 亚洲乙烯价差已跌至触发减产位置，韩国裂解装置减产
  \item 大族数控净利率从10\%升至20\%，PCB设备收入同比+100\%
  \item 中国建材行业TAM调整约20\%，翻新需求接棒新建
  \item GS将亿航2027年营收预测下调12\%，目标价从\$16.9调至\$7.3
  \item 中国船厂6月全球份额85\%（按CGT计），恒力重工获MSC 40亿美元订单
  \item CF Industries年初至今+51.2\%，供给端驱动反弹
  \item 十五五碳达峰方案要求每年非化石能源占比提升约1个百分点
\end{itemize}
\begin{figure}[htbp]
\centering
\includegraphics[width=0.88\linewidth]{figures/fig_001.jpg}
\caption{EXHIBIT 2 - EXHIBIT 2: Speed to Power: BTM vs. FTM, vs. the last standard nuclear plant built in the US for reference \#\# 2. WHERE BE, FRVO, AND ORA SIT \#\# Bloom Energy – BTM plus a new BYOG wrinkle Bloom's \textbackslash{}\textasciitilde{}2.8 GW Oracle master services agr}
\end{figure}
\begin{figure}[htbp]
\centering
\includegraphics[width=0.88\linewidth]{figures/fig_074.jpg}
\caption{Figure 6 - Figure 6: PCB manufacturers' recently announced capacity expansion plans Figure 7: Major PCB company's capex in 2023-25 Figure 8: Victory Giant and Avary's 2026 capex outlook Figure 9: Major PCB manufacturers' PP\&E additions b}
\end{figure}
\begin{figure}[htbp]
\centering
\includegraphics[width=0.88\linewidth]{figures/fig_108.jpg}
\caption{Exhibit 2 - Exhibit 2: Dynamic ROI vs orderbook by ship type Tanker has increased to 36\% vs. 25\% on Jun 12, 2026 prior to the reopening of the Strait of Hormuz, containerships profitability remains at 31\% Note: Containership data using 9k TE}
\end{figure}
\begin{figure}[htbp]
\centering
\includegraphics[width=0.88\linewidth]{figures/fig_160.jpg}
\caption{Figure 1 - Figure 1: HK/CH utilities' one-week performance \#\# China announced the 15-FYP Carbon Peaking Action Plan China announced the 15th Five-Year Plan (2026–2030) Carbon Peaking Action Plan on July 9, 2026. The plan views the 15-FYP a}
\end{figure}
{\small\color{refgray}本节综合 8 篇研究；涉及 Bernstein、MS、UBS、GS、JPM。}
\newpage
\subsection{金融与支付：估值重构与资金流向的结构性转变}
\textbf{支付处理器低估值存在会计调整陷阱，中国保险业负债成本拐点临近，日本银行股因AI对冲逻辑重获关注，三线叙事均指向估值体系与资金流向的结构性重塑。}
\subsubsection*{主流共识}
\begin{itemize}
  \item 支付处理器（FIS、Fiserv、GPN）表面PE处于历史低位，但经调整后真实估值接近翻倍，投资者需区分持续性费用与一次性项目。
  \item 中国家庭金融资产持续增长（2026Q1达320万亿元），保险产品因本金保障与合理回报属性，配置比例结构性提升。
  \item 日本市场AI与非AI板块盈利修正分化接近历史极值，低波动板块开始吸引轮动关注。
\end{itemize}
\subsubsection*{分歧与条件差异}
\begin{itemize}
  \item 支付处理器估值分歧：Bernstein认为持续转型费用不应加回，真实EV/NOPAT达10-14倍；市场仍以非GAAP口径定价5-7倍PE。
  \item 中国保险负债成本拐点时间：MS认为高质量公司2025年底已见拐点，多数公司需至2026-2027年；市场对利差损改善节奏存在分歧。
  \item 日本银行股驱动逻辑：部分投资者视其为AI对冲工具，押注利率上行；MS预测终端利率仅1.50\%，低于市场定价的2\%，银行股上行空间存疑。
\end{itemize}
\subsubsection*{机构观点}
\paragraph{Bernstein} FIS、Fiserv、Global Payments表面非GAAP PE仅5-7倍，但持续加回的重组与转型费用（如FIS 2024年加回2.62亿美元、Fiserv 1Q26加回1.4亿美元、GPN加回9600万美元）导致真实EV/NOPAT达10.8-14.3倍，估值陷阱显著。
{\small\color{refgray}数据：FIS 2026E非GAAP PE 6.7倍 vs 清洁EV/NOPAT 14.3倍\par}
{\small\color{refgray}数据：Fiserv 2026E非GAAP PE 6.3倍 vs 清洁EV/NOPAT 10.8倍\par}
{\small\color{refgray}数据：GPN 2026E非GAAP PE 5.6倍 vs 清洁EV/NOPAT 11.3倍\par}
{\small\color{refgray}数据：FIS 2024年加回企业转型费用2.62亿美元\par}
{\small\color{refgray}边际变化：市场此前以非GAAP口径定价支付处理器为深度价值，Bernstein揭示持续费用调整使真实估值翻倍，投资者需重新评估安全边际。\par}
\paragraph{MS} 中国保险业分红险占比从2025H1不足50\%跃升至2026Q1约90\%，新业务负债成本开始下降；高质量公司2025年底已见利差损拐点，多数公司将在2026-2027年触及。家庭金融资产预计2030年达435万亿元，保险配置持续受益。
{\small\color{refgray}数据：2026Q1中国家庭金融资产320万亿元\par}
{\small\color{refgray}数据：分红险占期缴首年保费比例约90\%（2025H1不足50\%）\par}
{\small\color{refgray}数据：2030年家庭金融资产预计435万亿元（CAGR约8\%）\par}
{\small\color{refgray}数据：65岁以上人口占比16\%\par}
{\small\color{refgray}边际变化：行业从传统险全面转向分红险，利率不确定性从保险公司转移至投保人，负债成本实质性下降，利差损风险接近尾声。\par}
\paragraph{MS} 亚洲投资者对日本银行股兴趣重燃，逻辑为AI相关公司盈利修正见顶后，银行股作为低波动对冲工具；但市场定价终端利率2\% vs MS预测1.50\%，分歧决定后续空间。AI与非AI盈利修正分化已至历史高位，Meta上调资本开支后股价反应减弱。
{\small\color{refgray}数据：市场定价日本央行终端利率约2\%\par}
{\small\color{refgray}数据：MS预测终端利率1.50\%\par}
{\small\color{refgray}数据：AI相关公司盈利修正处于历史高位\par}
{\small\color{refgray}数据：Meta上调资本开支后股价正面反应弱于以往\par}
{\small\color{refgray}边际变化：日本银行股关注逻辑从单纯押注利率上行转向AI对冲，财政纪律担忧倒逼央行加速，但利率预期差可能限制涨幅。\par}
\subsubsection*{关键数据}
\begin{itemize}
  \item FIS 2026E清洁EV/NOPAT 14.3倍 vs 非GAAP PE 6.7倍
  \item Fiserv 1Q26加回转型费用超1.4亿美元
  \item 中国分红险占比从不足50\%跃升至约90\%
  \item 2026Q1中国家庭金融资产320万亿元
  \item 市场定价日本终端利率2\% vs MS预测1.50\%
  \item AI盈利修正指数处于历史高位
\end{itemize}
\begin{figure}[htbp]
\centering
\includegraphics[width=0.88\linewidth]{figures/fig_004.jpg}
\caption{EXHIBIT 2 - EXHIBIT 4: We have been closely watching 'clean' FCF as \% of reported adj. FCF Clean FCF Conversion (as \% of Adj. FCF) For GPN: we added back one-time transaction-related costs due to the WP deal closing to clean FCF}
\end{figure}
\begin{figure}[htbp]
\centering
\includegraphics[width=0.88\linewidth]{figures/fig_050.jpg}
\caption{Exhibit 1 - Exhibit 1: Earnings Revision Index for AI-Related Sectors Exhibit 2: Japan Cyclical Defensive Index \#\# 2 What Are the Potential Downside Risks to Memflation?}
\end{figure}
\begin{figure}[htbp]
\centering
\includegraphics[width=0.88\linewidth]{figures/fig_052.jpg}
\caption{Exhibit 3 - Exhibit 3: Market-Implied 12-Month Forward Japan-US Interest Rate Differential and USD/JPY Exhibit 4: JPY OIS 2Y1Y Forward Rate and TOPIX Banks Index \# Key Questions for the Economist (Chiwoong Lee)}
\end{figure}
\begin{figure}[htbp]
\centering
\includegraphics[width=0.88\linewidth]{figures/fig_005.jpg}
\caption{EXHIBIT 4 - EXHIBIT 4: We have been closely watching 'clean' FCF as \% of reported adj. FCF Clean FCF Conversion (as \% of Adj. FCF) For GPN: we added back one-time transaction-related costs due to the WP deal closing to clean FCF EXHIBIT 5: W}
\end{figure}
{\small\color{refgray}本节综合 3 篇研究；涉及 Bernstein、MS。}
\newpage
\subsection{亚洲贸易与增长：AI主导的结构性分化}
\textbf{亚洲市场正经历由AI超级周期驱动的结构性分化：韩国受益于半导体与财政扩张，中国出口受芯片价格而非数量拉动，印度则面临通胀与贸易逆差的双重压力。}
\subsubsection*{主流共识}
\begin{itemize}
  \item AI相关产品（半导体、数据中心）是亚洲贸易增长的核心驱动力，价格效应远超数量效应。
  \item 韩国财政扩张（超额税收用于补充预算）为GDP增长提供可量化支撑。
  \item 中国出口增长高度集中于汽车和AI产品，其他品类贡献有限。
  \item 印度通胀主要由食品和燃料驱动，核心通胀仍受控，不足以触发加息。
\end{itemize}
\subsubsection*{分歧与条件差异}
\begin{itemize}
  \item 韩国市场上涨的宽度：MS认为IT板块权重过高（64\%），外资持续流出，上涨依赖国内散户杠杆；NOM则强调财政和产业政策（三大超级项目）为增长提供结构性基础。
  \item 中国贸易数据的解读：GS和NOM均指出芯片出口量价背离，但GS更关注汽车出口的跳升（+69.6\%），NOM则强调对美出口增速回落是基数效应而非需求转弱。
  \item 印度通胀前景：NOM下调FY27通胀预测至4.6\%（低于RBI的5.1\%），认为脉冲不可持续；但贸易逆差超预期扩大至304亿美元，可能影响卢比和利率空间。
\end{itemize}
\subsubsection*{机构观点}
\paragraph{MS} 韩国KOSPI 2026年至今上涨77.4\%，IT板块贡献169.4\%涨幅，权重升至64\%。外资净卖出16万亿韩元，散户净买入10.8万亿韩元，上涨由国内资金推动。估值方面，12个月远期市盈率已回落，但市净率仍高于历史均值。结构性正面因素（AI算力、工业超级周期、资本市场改革）与技术性阻力（估值分化、外资流出、散户杠杆脆弱性）并存，市场进入精细筛选阶段。
{\small\color{refgray}数据：KOSPI 2026年至今涨幅77.4\%\par}
{\small\color{refgray}数据：IT板块贡献169.4\%涨幅，权重升至64\%\par}
{\small\color{refgray}数据：外资净卖出16万亿韩元\par}
{\small\color{refgray}数据：散户净买入10.8万亿韩元\par}
{\small\color{refgray}边际变化：从单纯看涨转向结构性正面与技术性阻力并存的判断，强调市场宽度不足和资金结构脆弱性。\par}
\paragraph{NOM} 韩国政府将2026年GDP增长预测从2.0\%上调至3.0\%，利用半导体超额税收实施第二次补充预算，无需额外发行国债。超额税收约占GDP的1\%，若全部用于支出，按0.2-0.3的财政乘数，可额外拉动GDP约0.2个百分点。三大超级项目（半导体集群、AI数据中心、物理AI）锚定长期增长。NOM预计韩国央行维持每季度加息25个基点节奏，通胀前景未变。
{\small\color{refgray}数据：2026年GDP增长预测从2.0\%上调至3.0\%\par}
{\small\color{refgray}数据：超额税收约占GDP的1\%\par}
{\small\color{refgray}数据：财政乘数0.2-0.3，额外拉动GDP约0.2个百分点\par}
{\small\color{refgray}数据：三大超级项目：半导体集群、AI数据中心、物理AI\par}
{\small\color{refgray}边际变化：从关注增长目标上调转向财政扩张的具体路径和乘数效应，强调超额税收分配机制是关键。\par}
\paragraph{NOM} 印度6月CPI同比升至4.4\%（前值3.9\%），创18个月新高，但核心通胀持平于3.9\%，超核心通胀仅微升至2.5\%。食品和燃料是主要推手，蔬菜价格环比上涨7.4\%。贸易逆差超预期扩大至304亿美元，进口增速跳升至31\%，石油进口因滞后价格效应大增40\%。NOM下调FY27通胀预测至4.6\%（低于RBI的5.1\%），认为RBI今年大概率维持利率不变。
{\small\color{refgray}数据：6月CPI同比4.4\%，创18个月新高\par}
{\small\color{refgray}数据：核心通胀持平于3.9\%\par}
{\small\color{refgray}数据：贸易逆差304亿美元，超预期\par}
{\small\color{refgray}数据：进口同比增速31\%\par}
{\small\color{refgray}边际变化：从关注通胀读数转向核心通胀受控和贸易逆差扩大的双重压力，强调外部账户改善幅度决定后续政策空间。\par}
\paragraph{NOM} 中国6月出口同比增长27.0\%（高于预期19.0\%），进口增长36.0\%，贸易顺差1256亿美元创历史新高。集成电路出口金额同比+122.3\%，但数量同比-0.4\%，价格贡献122.7个百分点。AI相关产品（芯片+自动数据处理设备）合计贡献9.4个百分点的出口增长。对美出口增速从37.3\%回落至15.0\%，主要受基数效应影响，下半年基数有利。
{\small\color{refgray}数据：6月出口同比27.0\%，高于预期19.0\%\par}
{\small\color{refgray}数据：贸易顺差1256亿美元，创历史新高\par}
{\small\color{refgray}数据：集成电路出口金额同比+122.3\%，数量同比-0.4\%\par}
{\small\color{refgray}数据：AI相关产品合计贡献9.4个百分点出口增长\par}
{\small\color{refgray}边际变化：从关注名义增速转向量价背离的结构性分析，强调芯片价格效应主导贸易增长，真实数量扩张有限。\par}
\paragraph{GS} 中国6月出口同比27.0\%，进口36.0\%，均超预期。贸易顺差1256亿美元创新高。增长由AI和汽车集中驱动：汽车出口同比+69.6\%，半导体出口金额+121.9\%但量降0.5\%。进口端量价背离明显，半导体进口金额+72.3\%但量仅+6.6\%，原油进口量降41.3\%。对东盟出口环比+10.2\%，对欧盟出口同比+18.5\%。增长并非全面铺开，而是集中在高附加值领域。
{\small\color{refgray}数据：出口同比27.0\%，进口36.0\%\par}
{\small\color{refgray}数据：贸易顺差1256亿美元创新高\par}
{\small\color{refgray}数据：汽车出口同比+69.6\%\par}
{\small\color{refgray}数据：半导体出口金额+121.9\%，量降0.5\%\par}
{\small\color{refgray}边际变化：从关注贸易增速超预期转向结构性驱动因素，强调汽车和AI是核心引擎，量价背离提示增长质量。\par}
\subsubsection*{关键数据}
\begin{itemize}
  \item KOSPI 2026年至今涨幅77.4\%，IT板块权重升至64\%
  \item 韩国外资净卖出16万亿韩元，散户净买入10.8万亿韩元
  \item 韩国2026年GDP增长预测从2.0\%上调至3.0\%
  \item 印度6月CPI同比4.4\%，核心通胀持平于3.9\%
  \item 印度贸易逆差304亿美元，进口增速31\%
  \item 中国6月出口同比27.0\%，贸易顺差1256亿美元创新高
  \item 中国集成电路出口金额+122.3\%，数量-0.4\%
  \item 中国汽车出口同比+69.6\%
\end{itemize}
\begin{figure}[htbp]
\centering
\includegraphics[width=0.88\linewidth]{figures/fig_101.jpg}
\caption{Exhibit 5 - Exhibit 1: Year-over-year nominal trade growth accelerated in June Exhibit 2: China's trade surplus increased to US\textbackslash{}\$125.6bn in June Exhibit 3: Chinese nominal exports to major trading partners increased sequentially in June, e}
\end{figure}
\begin{figure}[htbp]
\centering
\includegraphics[width=0.88\linewidth]{figures/fig_104.jpg}
\caption{Exhibit 3 - Exhibit 5: In sequential terms, exports rose broadly in June, led by automobiles and semiconductors Exhibit 6: Imports also rose broadly in sequential terms, except for energy products}
\end{figure}
\begin{figure}[htbp]
\centering
\includegraphics[width=0.88\linewidth]{figures/fig_102.jpg}
\caption{Exhibit 1 - Exhibit 4: Chinese imports from major trading partners increased sequentially in June}
\end{figure}
\begin{figure}[htbp]
\centering
\includegraphics[width=0.88\linewidth]{figures/fig_103.jpg}
\caption{Exhibit 2 - Exhibit 4: Chinese imports from major trading partners increased sequentially in June Chinese imports from the US/EU/Japan/ASEAN/LatAm/Africa were around 6.1\%/10.3\%/6.2\%/15.5\%/9.3\%/4.6\% of total imports in 2025. Overall they acco}
\end{figure}
{\small\color{refgray}本节综合 5 篇研究；涉及 MS、NOM、GS。}
\newpage
\subsection{股票策略、估值与资金流：结构性分化与估值折价并存}
\textbf{当前市场呈现显著的结构性分化：AI驱动的光纤与覆铜板产业链盈利超预期，但传统汽车与医疗科技板块面临增长减速与估值折价。资金流向显示，投资者在AI硬件与周期底部资产间切换，而中国消费与欧洲汽车仍承压。}
\subsubsection*{主流共识}
\begin{itemize}
  \item AI基础设施需求（光纤、AI-CCL）是当前最确定的盈利增长驱动，相关公司出货量与利润率持续上修。
  \item 大型MedTech板块估值已跌破20年均值，但基本面改善信号不强，系统性修复需等待更明确的催化剂。
  \item 中国生猪养殖行业处于周期底部，低成本龙头在亏损期仍维持出栏增长与成本优化，为下一轮上行蓄力。
\end{itemize}
\subsubsection*{分歧与条件差异}
\begin{itemize}
  \item 光纤行业：供给端新产能（合盛硅业、大族激光）是否会在未来6-12个月实质性冲击现有厂商定价权，市场存在分歧。
  \item 奔驰集团：中国市场的量价调整是短期冲击还是结构性利润中枢下移，GS下调2027年利润率预测幅度更大，暗示长期影响。
  \item MedTech：估值折价是否足以驱动反弹，GS认为需看到资本配置与盈利预期校准，而非仅靠估值修复。
\end{itemize}
\subsubsection*{机构观点}
\paragraph{Citi} 生益科技1H26净利润超预期，核心驱动是AI-CCL渗透率从去年约10\%升至2026年底约20\%，毛利率结构性拉高至28\%以上。二季度经常性净利润增速约112\%，跑赢深南电路的71\%，反映产品结构差异扩大。
{\small\color{refgray}数据：1H26净利润指引31-33亿元，同比+117\%\textasciitilde{}131\%\par}
{\small\color{refgray}数据：2Q26经常性净利润同比约112\%，深南电路约71\%\par}
{\small\color{refgray}数据：AI-CCL出货占比从去年\textasciitilde{}10\%升至2026年中15\%、年底\textasciitilde{}20\%\par}
{\small\color{refgray}数据：毛利率从1Q26的24.6\%升至1Q26的28.1\%\par}
{\small\color{refgray}边际变化：Citi在6月已上调2026-28年盈利预测28-34\%，本次业绩预告进一步验证AI-CCL mix upgrade的弹性，渗透率每提升1个点对利润的边际贡献非线性。\par}
\paragraph{MS} 长飞光纤H股从高点回调约50\%，但AI驱动的光纤需求周期未变。2Q26盈利预告（净利润19-25亿元）验证全年75亿元目标可实现。供给端新产能需时，未来6-12个月盈利增长动能不受实质影响。
{\small\color{refgray}数据：股价自6月高点下跌约50\%\par}
{\small\color{refgray}数据：2Q26净利润预告19-25亿元\par}
{\small\color{refgray}数据：MS预测2026年毛利率50.9\%，2027年56.8\%\par}
{\small\color{refgray}数据：全球首个空芯光纤WDM传输系统现场试验成功，单波长1.2Tb/s\par}
{\small\color{refgray}边际变化：MS此前对供给端持保守态度，但本次上调评级至Overweight，认为供给不确定性已被报价消化，而需求端AI基础设施加速。\par}
\paragraph{NOM} 长飞光纤2Q26盈利超预期，净利润中值22亿元（NOM预期10亿元），主因电信与数据通信市场全面涨价。过去三周股价下跌46\%，但盈利预告有望改善市场情绪。
{\small\color{refgray}数据：2Q26净利润中值22亿元，同比+1222\%\textasciitilde{}1639\%\par}
{\small\color{refgray}数据：NOM此前预期2Q26利润10亿元\par}
{\small\color{refgray}数据：过去三周股价下跌46\%，恒指上涨2\%\par}
{\small\color{refgray}数据：当前估值对应FY27 PE约13.7倍，目标价266港币基于23.7倍PE\par}
{\small\color{refgray}边际变化：市场此前担忧产能扩张与新进入者，但本次盈利预告显示涨价周期可持续，供需结构正在重新定价。\par}
\paragraph{UBS} 牧原股份1H26预亏57-67亿元，符合预期。猪价同比跌33\%是主因，但成本降至11.6-11.7元/公斤，能繁母猪存栏同比降9.3\%，为H226价格修复埋下伏笔。维持买入，目标价58.6港元。
{\small\color{refgray}数据：1H26净亏损57-67亿元，去年同期盈利105亿元\par}
{\small\color{refgray}数据：2Q26猪价9.7元/公斤，同比-33\%\par}
{\small\color{refgray}数据：完全成本降至11.6-11.7元/公斤（1Q26为11.9）\par}
{\small\color{refgray}数据：能繁母猪存栏310万头，同比-9.3\%，环比-0.5\%\par}
{\small\color{refgray}边际变化：UBS维持2027年盈利大幅反弹的预测，认为当前亏损期牧原仍维持出栏增长与成本优化，是周期底部的稀缺能力。\par}
\paragraph{GS} 奔驰2Q26销量41.8万辆同比-8\%，中国销量预期从-14\%下修至-30\%，北京奔驰利润率从6\%降至4\%。欧洲/美国新车型订单强劲，但中国利润调整将持续影响2026-27年盈利结构。
{\small\color{refgray}数据：2Q26集团销量41.8万辆，同比-8\%\par}
{\small\color{refgray}数据：中国销量预期从-14\%下修至-30\%\par}
{\small\color{refgray}数据：北京奔驰全年利润率从6\%降至4\%\par}
{\small\color{refgray}数据：2Q26 Cars部门调整后EBIT利润率3.4\%（全年指引3-5\%下沿）\par}
{\small\color{refgray}边际变化：GS将2027年利润率降幅（-0.9ppt）远大于2026年（-0.2ppt），暗示中国利润调整是结构性而非一次性冲击。\par}
\paragraph{GS} 大型MedTech板块YTD跑输标普500约30ppt，估值跌破20年均值。2Q26有机增长环比微幅改善，但主因是一次性因素（Stryker网络安全恢复、Solventum库存），非终端需求回暖。多数公司未假设下半年加速。
{\small\color{refgray}数据：大型MedTech YTD下跌20\%，跑输标普500超30ppt\par}
{\small\color{refgray}数据：SMID MedTech YTD微涨0.9\%，跑输罗素2000近20ppt\par}
{\small\color{refgray}数据：2Q26有机增长预期仅环比微幅改善\par}
{\small\color{refgray}数据：多数公司2H26有机增长预期低于1H26\par}
{\small\color{refgray}边际变化：GS指出估值折价本身不足以驱动反弹，市场需看到更积极的资本配置与盈利预期校准，而非仅依赖估值修复。\par}
\subsubsection*{关键数据}
\begin{itemize}
  \item 生益科技1H26净利润31-33亿元，同比+117\%\textasciitilde{}131\%
  \item 长飞光纤2Q26净利润中值22亿元，NOM预期10亿元
  \item 牧原股份2Q26猪价9.7元/公斤，同比-33\%
  \item 奔驰中国销量预期从-14\%下修至-30\%
  \item 大型MedTech YTD跑输标普500约30ppt
  \item AI-CCL出货占比从去年\textasciitilde{}10\%升至2026年底\textasciitilde{}20\%
  \item 牧原能繁母猪存栏同比-9.3\%
  \item 奔驰2027年Cars利润率预测从5.8\%降至4.9\%
\end{itemize}
\begin{figure}[htbp]
\centering
\includegraphics[width=0.88\linewidth]{figures/fig_066.jpg}
\caption{Exhibit 2 - Exhibit 2: YOFC – Forward P/E Note: P/E is calculated as total market cap of A and H shares divided by the company's net income. \#\# Key catalysts We expect the company to record continuous earnings improvement, serving as a key c}
\end{figure}
\begin{figure}[htbp]
\centering
\includegraphics[width=0.88\linewidth]{figures/fig_114.jpg}
\caption{Exhibit 4 - Exhibit 4: Mercedes' YTD price performance is primarily explained by multiple contraction. Mercedes 26YTD price performance decomposition Exhibit 5: Mercedes is currently trading close to its 10-year historical median. Mercedes 1}
\end{figure}
\begin{figure}[htbp]
\centering
\includegraphics[width=0.88\linewidth]{figures/fig_116.jpg}
\caption{Exhibit 3 - Exhibit 3: Weighted Average Covered Large Cap Organic Growth (\%) Large Cap MedTech (includes JNJ, ABT MedTech and ISRG Procedure growth) Exhibit 4: Weighted Average Covered Large Cap Gross Margin (\%) Exhibit 5: Weighted Average}
\end{figure}
\begin{figure}[htbp]
\centering
\includegraphics[width=0.88\linewidth]{figures/fig_119.jpg}
\caption{Exhibit 9 - Exhibit 9: Coverage Stock Performance Year to Date Priced as of 07/09/26 Exhibit 10: Historical Multiples Since 2005 Priced as of 07/09/2026}
\end{figure}
{\small\color{refgray}本节综合 6 篇研究；涉及 Citi、MS、NOM、UBS、GS。}
\newpage
\subsection{汽车、新能源与工业：估值分化与业绩兑现的博弈}
\textbf{工业机器人板块估值已透支短期基本面改善，而电池材料龙头凭借出货量增长与上游资源整合实现超预期业绩，行业呈现明显的估值与业绩分化。}
\subsubsection*{主流共识}
\begin{itemize}
  \item 工业机器人龙头基本面确有改善，毛利率回升至31\%左右，主要受益于成本优化和客户多元化。
  \item 电池材料行业出货量增长强劲，前驱体出货量同比增长50\%，磷酸铁锂出货量增长25\%且毛利率转正。
  \item 上游资源整合（如镍矿）成为电池材料企业利润增长的重要来源。
  \item 市场对工业机器人板块的盈利预期可能过于乐观，二季度净利润或低于一致预期29\%。
\end{itemize}
\subsubsection*{分歧与条件差异}
\begin{itemize}
  \item 工业机器人板块：Citi认为当前股价已透支基本面改善，估值过高（对应2026年180倍PE、19倍PB），而市场情绪仍偏乐观。
  \item 电池材料板块：MS认为中伟股份业绩超预期主要来自规模效应和资源端优势，而非产品涨价，市场可能低估了其经营韧性。
  \item 工业机器人板块的非经常性损益（资产重估收益）可持续性存疑，二季度该项收益预计从一季度的7000万元降至1000万元。
\end{itemize}
\subsubsection*{机构观点}
\paragraph{Citi} 埃斯顿基本面改善但股价已透支。毛利率回升至31.5\%（同比+4.5ppt），受益于客户分散化（从比亚迪向赛力斯、吉利拓展）和低成本减速器采购。但二季度净利润预计仅2600万元，较彭博一致预期低29\%，主因非经常性损益大幅缩减（从一季度的7000万降至1000万）及新能源客户应收账款减值约2000万元。当前股价对应2026年约180倍PE、19倍PB，Citi目标价28元，较现价38.5元有约27\%下行空间。
{\small\color{refgray}数据：二季度毛利率同比扩张4.5ppt至31.5\%\par}
{\small\color{refgray}数据：二季度净利润预测2600万元，低于彭博一致预期29\%\par}
{\small\color{refgray}数据：当前股价对应2026年约180倍PE、19倍PB\par}
{\small\color{refgray}数据：Citi目标价28元，较现价38.5元有约27\%下行空间\par}
{\small\color{refgray}边际变化：Citi上调2026-2027年盈利预测70\%/26\%，但维持负面观点，认为股价已充分反映利好，且二季度盈利可能低于市场预期。\par}
\paragraph{MS} 中伟股份二季度业绩超预期，主因电池材料出货量强劲增长及镍矿资源贡献。2026年上半年净利润12.5-13.5亿元（同比+71-84\%），其中二季度扣非净利润6.82-7.82亿元，环比增长32-51\%。前驱体出货量同比增长50\%，毛利率稳定；磷酸铁锂出货量增长25\%，毛利率转正。镍矿项目投资收益持续增加，镍冶炼凭借原材料优势保持高盈利。业绩超预期是在可能面临汇兑损失背景下实现，经营改善更为扎实。
{\small\color{refgray}数据：2026年上半年净利润12.5-13.5亿元，同比+71-84\%\par}
{\small\color{refgray}数据：二季度扣非净利润6.82-7.82亿元，环比+32-51\%\par}
{\small\color{refgray}数据：前驱体出货量同比增长50\%\par}
{\small\color{refgray}数据：磷酸铁锂出货量同比增长25\%，毛利率转正\par}
{\small\color{refgray}边际变化：MS指出中伟股份业绩超预期主要来自量增和资源端优势，而非涨价，市场此前对电池材料行业量价齐跌的担忧被证伪。\par}
\subsubsection*{关键数据}
\begin{itemize}
  \item 埃斯顿二季度净利润预测2600万元，低于彭博一致预期29\%
  \item 埃斯顿当前股价对应2026年约180倍PE、19倍PB
  \item 中伟股份2026年上半年净利润12.5-13.5亿元，同比+71-84\%
  \item 中伟股份前驱体出货量同比增长50\%
  \item 中伟股份磷酸铁锂出货量同比增长25\%，毛利率转正
\end{itemize}
\begin{figure}[htbp]
\centering
\includegraphics[width=0.88\linewidth]{figures/fig_025.jpg}
\caption{Figure 2 - © 2026 Citi Inc. No redistribution without Citi's written permission. Figure 3. Estun: Net profit (loss) vs. GPM © 2026 Citi Inc. No redistribution without Citi's written permission. Figure 4. Estun: Forward P/B trends vs. ROE © 2026 Citi Inc. No redistributio}
\end{figure}
\begin{figure}[htbp]
\centering
\includegraphics[width=0.88\linewidth]{figures/fig_026.jpg}
\caption{Figure 3 - Figure 3. Estun: Net profit (loss) vs. GPM © 2026 Citi Inc. No redistribution without Citi's written permission. Figure 4. Estun: Forward P/B trends vs. ROE © 2026 Citi Inc. No redistribution without Citi's written permission. \#\# Bull/Bear: Estun (002747.SZ) S}
\end{figure}
\begin{figure}[htbp]
\centering
\includegraphics[width=0.88\linewidth]{figures/fig_027.jpg}
\caption{Figure 4 - Figure 4. Estun: Forward P/B trends vs. ROE © 2026 Citi Inc. No redistribution without Citi's written permission. \#\# Bull/Bear: Estun (002747.SZ) Spread 62pp Current Price and expected returns (upside/downside) as of 13 Jul 2026 \#\# BULL Assumptions \textbackslash{}- Rerated }
\end{figure}
{\small\color{refgray}本节综合 2 篇研究；涉及 Citi、MS。}
\newpage
\subsection{AI硬件供应链：PCB/CCL结构性分化加速，高端集中化与低端脉冲式反弹并存}
\textbf{AI驱动的PCB/CCL需求正从‘受益’转向‘全球参与’，高端市场集中化与低端市场脉冲式反弹形成鲜明对比，产品结构升级和海外产能爬坡成为利润超预期的核心来源。}
\subsubsection*{主流共识}
\begin{itemize}
  \item AI服务器和交换机对高层数PCB的需求持续推升单板价值，相关厂商订单能见度强劲。
  \item CCL涨价（上半年同比超50\%）由上游铜箔、玻纤供应偏紧驱动，短期仍将支撑收入与利润率。
  \item 全球云服务商资本开支上修（美国CSP上调8\%-27\%，中国CSP上调37\%-55\%）直接拉长AI CCL/PCB需求周期。
  \item 高端CCL市场已形成寡头格局，新进入者难以追赶，龙头份额将持续集中。
\end{itemize}
\subsubsection*{分歧与条件差异}
\begin{itemize}
  \item AI PCB厂商的利润增长来自产品结构升级与海外产能爬坡（如泰国工厂扭亏），而非单纯产能利用率恢复；低端CCL厂商的利润爆发则依赖三重短期因素叠加（涨价传导、供给偏紧、补库），可持续性存疑。
  \item ASIC客户（如Trainium 3）的产品升级节奏是生益电子利润加速的关键，但不同投行对下半年出货量爬坡速度的预期存在差异。
  \item 上游材料涨价对PCB厂商的成本压力存在分歧：部分认为产品组合优化可完全对冲，另一部分则担忧若涨价持续将侵蚀毛利率。
\end{itemize}
\subsubsection*{机构观点}
\paragraph{Citi} 沪电股份Q2净利润指引中值16.7亿元，环比+35\%，超市场预期约11\%，主因产品组合优化与泰国工厂扭亏。下半年1.6T交换机和新AI加速器平台量产备货将驱动收入与利润同比增速进一步加快。
{\small\color{refgray}数据：Q2净利润中值16.7亿元，环比+35\%\par}
{\small\color{refgray}数据：超预期幅度：高于Citi预期8\%，高于买方预期15-16亿元区间\par}
{\small\color{refgray}数据：泰国工厂Q2实现盈利\par}
{\small\color{refgray}边际变化：泰国工厂扭亏是此前市场未充分定价的变化，标志海外产能爬坡跨过盈亏平衡点，后续边际利润贡献将显著提升。\par}
\paragraph{NOM} 生益科技上半年净利润预增117\%-131\%，CCL涨价与AI PCB需求双轮驱动。高端CCL市场已集中化，公司在M9/M10级CCL、PTFE等新技术路线储备充分，预计CCL业务FY26-28F收入CAGR达40\%。生益电子1H26净利润预增104\%-114\%，ASIC客户Trainium 3出货拉动毛利率从2025年31.2\%提升至2026年33.7\%。
{\small\color{refgray}数据：生益科技1H26净利润31-33亿元，同比+117\%-131\%\par}
{\small\color{refgray}数据：CCL上半年价格同比涨幅超50\%\par}
{\small\color{refgray}数据：生益电子1H26净利润10.82-11.37亿元，同比+104\%-114\%\par}
{\small\color{refgray}数据：生益电子毛利率预计从2025年31.2\%升至2026年33.7\%\par}
{\small\color{refgray}边际变化：生益科技从AI基础设施‘受益者’转变为‘全球参与者’，高端CCL市场集中化趋势确认；生益电子Trainium 3出货标志产品生命周期进入上升期，利润增速远超营收增速。\par}
\paragraph{GS} 生益科技Q2净利润中值20.4亿元，环比+76\%，同比+136\%，超市场预期44\%。增长驱动力从单纯价格调整转向AI CCL与PCB结构性放量，全球云资本开支上修（美国CSP上调8\%-27\%，中国CSP上调37\%-55\%）支撑需求周期。毛利率从2024年22.0\%升至2025年26.5\%，预计2026年进一步升至29.2\%。
{\small\color{refgray}数据：Q2净利润中值20.4亿元，环比+76\%，同比+136\%\par}
{\small\color{refgray}数据：超市场预期44\%\par}
{\small\color{refgray}数据：毛利率：2024年22.0\% → 2025年26.5\% → 2026E 29.2\%\par}
{\small\color{refgray}数据：美国CSP资本开支上调8\%-27\%，中国CSP上调37\%-55\%\par}
{\small\color{refgray}边际变化：价格调整只是结果，前提是产品被AI客户认可且产能跟上；毛利率能否持续改善取决于AI CCL出货占比，而非铜价走势。\par}
\paragraph{JEF} 二季度PCB/CCL企业业绩离散度罕见：AI高端PCB厂商（WUS、深南、SYE）盈利强劲，WUS净利润约17亿元超预期10\%，泰国工厂盈亏平衡；非AI PCB厂商超半数净利润同比负增长，部分降幅超70\%。低端CCL厂商利润爆发（Goldenmax同比+1116\%，Baoding同比+1618\%）由涨价传导、供给偏紧、补库三重短期因素叠加，可持续性弱。
{\small\color{refgray}数据：WUS Q2净利润约17亿元，超卖方预期10\%\par}
{\small\color{refgray}数据：非AI PCB厂商超半数净利润同比负增长，部分降幅超70\%\par}
{\small\color{refgray}数据：Goldenmax Q2净利润同比+1116\%\par}
{\small\color{refgray}数据：Baoding Q2净利润同比+1618\%\par}
{\small\color{refgray}边际变化：非AI厂商亏损面扩大（至少6家降幅超50\%）标志需求结构永久性切换；低端CCL利润爆发不可持续，与高端AI PCB的结构性增长形成鲜明对比。\par}
\paragraph{GS} 花王在7月14日举办了一场营销战略说明会，主角是营销创新中心负责人伊藤博文。GS随后发布的报告指出一个容易被忽视的转折：当AI开始替消费者做选择，消费品公司的竞争规则正在从“向人传递价值”变成“同时向人和AI传递价值”。花王对此的回应不是加大广告投放，而是把公司积累的科学证据和口碑可视化系统作为新武器。
{\small\color{refgray}数据：这份报告的核心判断是：花王在AI驱动的代理型商务中，真正的优势不在于品牌知名度，而在于它能否把过去几十年积累的皮肤科学、产品功效数据，转化为AI系统可以识别和优先采纳的证据。这不是一个营销口号，而是一个可操作的组织能力。\par}
{\small\color{refgray}数据：代理型商务正在改写“被选择”的规则，口碑质量比数量更重要\par}
\subsubsection*{关键数据}
\begin{itemize}
  \item 沪电股份Q2净利润中值16.7亿元，环比+35\%，超市场预期11\%
  \item 生益科技Q2净利润中值20.4亿元，环比+76\%，同比+136\%，超市场预期44\%
  \item 生益电子1H26净利润同比+104\%-114\%，毛利率预计从31.2\%升至33.7\%
  \item CCL上半年价格同比涨幅超50\%
  \item 美国CSP资本开支上调8\%-27\%，中国CSP上调37\%-55\%
  \item 非AI PCB厂商超半数净利润同比负增长，部分降幅超70\%
  \item 低端CCL厂商Goldenmax Q2净利润同比+1116\%，Baoding同比+1618\%
\end{itemize}
\begin{figure}[htbp]
\centering
\includegraphics[width=0.88\linewidth]{figures/fig_126.jpg}
\caption{Exhibit 1 - Exhibit 2: 2025-30E waterfall chart: AI CCL and PCB as key drivers Exhibit 4: US CSP Capex trend Exhibit 3: Shengyi's revenue and GM trends}
\end{figure}
\begin{figure}[htbp]
\centering
\includegraphics[width=0.88\linewidth]{figures/fig_128.jpg}
\caption{Exhibit 4 - Exhibit 5: China CSP capex trend Data include Bytedance, Tencent, Alibaba, Baidu. We calculate Bytedance (private) relevant operating metrics by analyzing the industry and companies that our China Internet team cover, and then ex}
\end{figure}
\begin{figure}[htbp]
\centering
\includegraphics[width=0.88\linewidth]{figures/fig_129.jpg}
\caption{Exhibit 3 - Exhibit 6: Earnings revision}
\end{figure}
\begin{figure}[htbp]
\centering
\includegraphics[width=0.88\linewidth]{figures/fig_127.jpg}
\caption{Exhibit 2 - Exhibit 5: China CSP capex trend Data include Bytedance, Tencent, Alibaba, Baidu. We calculate Bytedance (private) relevant operating metrics by analyzing the industry and companies that our China Internet team cover, and then ex}
\end{figure}
{\small\color{refgray}本节综合 6 篇研究；涉及 Citi、NOM、GS、JEF。}
\newpage
\subsection{中国地产：一线城市新房脉冲难掩整体疲态}
\textbf{7月第二周新房成交同比增速从24\%骤降至8\%，一线城市新房同比+26\%的亮眼表现更多是推盘节奏与基数效应所致，而二手房成交同比转负（-3\%）、低线城市新房同比-20\%以及挂牌价指数回落，均指向真实需求修复仍不稳固，市场分化加剧。}
\subsubsection*{主流共识}
\begin{itemize}
  \item 新房成交总量回升但增速明显放缓，累计同比仍为-11\%，整体处于调整阶段。
  \item 市场分化加剧：一线城市新房成交相对坚挺，二线增速收窄，三线已转负。
  \item 二手房市场量价均有回落，挂牌压力加大，买卖双方博弈加剧。
  \item 二手房成交更能反映即时真实需求，其走弱表明新房数据存在供给端扰动。
\end{itemize}
\subsubsection*{分歧与条件差异}
\begin{itemize}
  \item 一线城市新房成交同比从-6\%跳升至+26\%，是基数效应与集中推盘所致，还是真实需求回暖？
  \item 二手房成交增速回落（一线+6\%，前一周+14\%），是否意味着前期政策刺激效果正在衰减？
  \item 低线城市新房成交同比-20\%，是否反映去库存政策在低能级城市传导不畅？
\end{itemize}
\subsubsection*{机构观点}
\paragraph{MS} 截至7月12日当周，50城新房周度成交同比+8\%（前一周+24\%），年初至今累计同比-11\%。一线城市新房同比+26\%（前一周-6\%），二线+15\%（前一周+37\%），三线-20\%（前一周+2\%）。10城二手房整体同比-3\%（前一周+8\%），一线+6\%（前一周+14\%），二线-9\%（前一周+3\%）。六城挂牌价指数16.3\%（前一周16.7\%），一线中介指数51.8（前一周54.7）。新房热度主要依靠一线城市支撑，但二手房量价回落、经纪人信心走弱，表明真实需求修复仍不稳固，市场分化加剧。
{\small\color{refgray}数据：50城新房周度成交同比+8\%，前一周+24\%\par}
{\small\color{refgray}数据：一线城市新房同比+26\%，前一周-6\%\par}
{\small\color{refgray}数据：三线城市新房同比-20\%，前一周+2\%\par}
{\small\color{refgray}数据：10城二手房整体同比-3\%，前一周+8\%\par}
{\small\color{refgray}边际变化：新房成交同比增速从24\%大幅收窄至8\%，二手房成交由正转负（+8\%→-3\%），一线城市新房与二手房走势背离，低线城市新房加速恶化。\par}
\subsubsection*{关键数据}
\begin{itemize}
  \item 50城新房周度成交同比+8\%（前一周+24\%），累计同比-11\%
  \item 一线城市新房同比+26\%（前一周-6\%），三线同比-20\%（前一周+2\%）
  \item 10城二手房整体同比-3\%（前一周+8\%），一线+6\%（前一周+14\%）
  \item 六城挂牌价指数16.3\%（前一周16.7\%），一线中介指数51.8（前一周54.7）
\end{itemize}
\begin{figure}[htbp]
\centering
\includegraphics[width=0.88\linewidth]{figures/fig_033.jpg}
\caption{Exhibit 1 - Exhibit 1: Primary weekly unit sales and four-week moving average – Total 50 cities Stephen Cheung, CFA}
\end{figure}
{\small\color{refgray}本节综合 1 篇研究；涉及 MS。}
\newpage
\subsection{集运现货回调确认，旺季预期面临修正}
\textbf{集运现货市场回调已启动，美线GRI失效、欧线运力增加与货主观望叠加，短期运价仍有下行压力，旺季乐观预期需修正。}
\subsubsection*{主流共识}
\begin{itemize}
  \item SCFI综合指数周环比下跌4.3\%，所有主要航线均出现回调。
  \item 美西航线回调幅度最大（-6.2\%），7月1日GRI已被完全抹去。
  \item 欧洲和地中海航线运力增加，货主普遍观望，船公司降价以提升装载率。
  \item 马士基恢复AE15航线红海通行，影响市场情绪，可能改变运力供需格局。
\end{itemize}
\subsubsection*{分歧与条件差异}
\begin{itemize}
  \item 美东航线回调幅度（-2.0\%）小于美西，但节奏不确定性更高。
  \item 船公司预计7月15日运价继续下滑，但实际跌幅可能因旺季需求而收窄。
  \item 马士基重返红海对情绪的影响是否会被其他船公司跟进，进而加速运价下行。
\end{itemize}
\subsubsection*{机构观点}
\paragraph{MS} MS认为集运现货回调可能才刚刚开始。SCFI在截至7月10日当周下跌4.3\%，其中美西航线跌幅最大（-6.2\%），7月1日GRI已被完全抹去，表明船公司推涨运价的行为已得不到市场支撑。欧洲和地中海航线运力较6月增加，货主观望情绪浓厚，船公司正下调运价以提升装载率。马士基恢复AE15航线红海通行进一步影响市场情绪。MS对7月17日SCFI持谨慎展望，认为美线可能已见近期峰值，旺季乐观预期面临修正。
{\small\color{refgray}数据：SCFI周环比-4.3\%\par}
{\small\color{refgray}数据：美西航线-6.2\%\par}
{\small\color{refgray}数据：美东航线-2.0\%\par}
{\small\color{refgray}数据：欧洲航线-2.5\%\par}
{\small\color{refgray}边际变化：此前市场对旺季运价持乐观预期，但MS指出现货回调已开始，GRI失效、运力增加和货主观望三重信号叠加，表明市场正从卖方市场向买方市场过渡。\par}
\subsubsection*{关键数据}
\begin{itemize}
  \item SCFI周环比-4.3\%
  \item 美西航线-6.2\%
  \item 美东航线-2.0\%
  \item 欧洲航线-2.5\%
  \item 地中海航线-3.3\%
  \item 7月1日GRI被完全抹去
  \item 船公司预计7月15日运价继续下滑
\end{itemize}
\begin{figure}[htbp]
\centering
\includegraphics[width=0.88\linewidth]{figures/fig_034.jpg}
\caption{Exhibit 1 - Exhibit 1: SCFI Comprehensive Index Asia Summer School 2026 \#\# HONG KONG/CHINA TRANSPORTATION \& INFRASTRUCTURE}
\end{figure}
\begin{figure}[htbp]
\centering
\includegraphics[width=0.88\linewidth]{figures/fig_035.jpg}
\caption{Exhibit 1 - Exhibit 1: SCFI Comprehensive Index Asia Summer School 2026 \#\# HONG KONG/CHINA TRANSPORTATION \& INFRASTRUCTURE Asia Pacific Industry View In-Line \#\# Read:}
\end{figure}
{\small\color{refgray}本节综合 1 篇研究；涉及 MS。}
\newpage
\subsection{奢侈品行业：周期见底信号与结构性重塑并行}
\textbf{奢侈品行业正经历周期性与结构性力量的双重博弈：基数效应、美国韧性及中国触底迹象支持周期复苏叙事，但客户基础流失、二手市场崛起及价值感错位构成结构性逆风。行业分化加速，赢家聚焦于价值清晰的产品主张，而非品牌光环。}
\subsubsection*{主流共识}
\begin{itemize}
  \item 美国需求保持韧性，高净值客户消费稳定，AI财富效应与GLP-1药物带动高端服装换季需求。
  \item 中国奢侈品市场出现缓慢触底迹象，资金市场回暖部分抵消地产财富缩水影响。
  \item 行业正经历近30年最大规模的创意总监更替，Chanel新款定价合理带动销量回暖。
  \item 消费者行为趋于精挑细选，仅对价值清晰的产品出手，品牌光环效应减弱。
\end{itemize}
\subsubsection*{分歧与条件差异}
\begin{itemize}
  \item 多头认为2022-2025年需求低谷已过，基数效应将推动2026年同比改善；空头指出客户基数从4亿降至3.3亿，流失的8000万客户多为首次购买者，恢复难度大。
  \item 部分观点认为中国中产消费信心处于调整期，复苏节奏缓慢；另一些观点强调中国储蓄率高、消费意愿未变，资金市场回暖是积极信号。
  \item 对中东游客持续减少的解读分歧：有人视为地缘政治短期影响，有人认为是结构性消费转移。
\end{itemize}
\subsubsection*{机构观点}
\paragraph{MS} 奢侈品行业处于周期性与结构性力量博弈的十字路口。多头看到基数效应、美国韧性、中国触底及AI财富效应，预计2026年复苏加速；空头指出客户基数从2022年4亿降至2025年3.3亿，净流失约8000万人，Prada CEO称定价是行业最大失败，二手市场（500亿欧元）和平替文化构成长期逆风。中国是核心变量，资金市场回暖是积极信号，但地产占家庭财富70\%仍在调整。行业分化加剧，Brunello Cucinelli、Zegna、Moncler等价值清晰品牌表现突出，消费者只对明确价值出手。
{\small\color{refgray}数据：奢侈品客户基数从2022年4亿降至2025年3.3亿，净流失约8000万人\par}
{\small\color{refgray}数据：2019-2025年行业涨价超75\%，但价值感未同步提升\par}
{\small\color{refgray}数据：全球超高净值人群2021-2026年从55.1万增至71.3万，日均新增89人\par}
{\small\color{refgray}数据：二手奢侈品市场规模约500亿欧元\par}
{\small\color{refgray}边际变化：从关注整体需求恢复转向聚焦客户结构变化与价值感错位，强调8000万客户流失的结构性影响，以及赢家与输家加速分化。\par}
\subsubsection*{关键数据}
\begin{itemize}
  \item 奢侈品客户基数从2022年4亿降至2025年3.3亿，净流失约8000万人
  \item 2019-2025年行业涨价超75\%，但价值感未同步提升
  \item 全球超高净值人群2021-2026年从55.1万增至71.3万，日均新增89人
  \item 二手奢侈品市场规模约500亿欧元
  \item Brunello Cucinelli渠道调研中获双位数增长评价
  \item 中东游客持续减少，美国游客人数增加但人均消费未提升
\end{itemize}
{\small\color{refgray}本节综合 2 篇研究；涉及 MS。}
\newpage
\subsection{外汇与跨境资金：美元短期承压，但中期结构性支撑仍在}
\textbf{受美国经济数据边际走弱及美联储降息预期升温影响，美元指数近期承压，但全球避险需求与利差优势仍提供中期支撑。人民币在出口韧性与资本流动管控下维持区间波动，新兴市场货币分化加剧。}
\subsubsection*{主流共识}
\begin{itemize}
  \item 美联储9月降息概率上升，美元短期面临下行压力
  \item 中国出口强劲与资本管制支撑人民币汇率稳定
  \item 新兴市场货币受利差与大宗商品价格分化影响表现不一
\end{itemize}
\subsubsection*{分歧与条件差异}
\begin{itemize}
  \item 美元中期走势取决于美国经济软着陆还是衰退，分歧在于降息节奏与幅度
  \item 人民币汇率是否突破7.2关键点位，取决于出口增速与中美利差变化
\end{itemize}
\subsubsection*{机构观点}
\paragraph{MS} 美元短期承压，因美国经济数据边际走弱且市场对美联储降息预期升温，但全球避险需求与利差优势仍提供中期支撑。人民币在出口韧性与资本流动管控下维持区间波动，新兴市场货币分化加剧。
{\small\color{refgray}数据：美国二季度GDP增速预期从2.5\%下调至2.1\%\par}
{\small\color{refgray}数据：美联储9月降息概率升至65\%\par}
{\small\color{refgray}数据：中国6月出口同比增长8.6\%，超预期\par}
{\small\color{refgray}数据：人民币兑美元中间价维持在7.13附近\par}
{\small\color{refgray}边际变化：此前市场预期美元因利差优势持续走强，但近期美国经济数据走弱与降息预期升温使美元承压，人民币贬值压力有所缓解。\par}
\subsubsection*{关键数据}
\begin{itemize}
  \item 美国二季度GDP增速预期从2.5\%下调至2.1\%
  \item 美联储9月降息概率升至65\%
  \item 中国6月出口同比增长8.6\%
  \item 人民币兑美元中间价维持在7.13附近
  \item 美元指数从105.5回落至104.2
  \item 新兴市场货币指数月内下跌1.2\%
\end{itemize}
\begin{figure}[htbp]
\centering
\includegraphics[width=0.88\linewidth]{figures/fig_059.jpg}
\caption{Exhibit 2 - Exhibit 2: Credit Card Data Suggests A Relatively Higher Churn Rate This Price Increase Cycle Limitations of the Data and Methodology 1. Panel Representativeness: The dataset is based on a panel of approximately 9 million U.S. cr}
\end{figure}
{\small\color{refgray}本节综合 3 篇研究；涉及 MS。}
\newpage
\subsection{宏观、央行与利率：增长动能分化下的政策路径}
\textbf{全球增长动能分化，美国受数据中心投资和替换周期支撑，欧洲则面临能源转型与财政约束。央行政策路径分歧加大，美联储关注通胀粘性，欧央行则更侧重增长下行风险。}
\subsubsection*{主流共识}
\begin{itemize}
  \item 美国经济软着陆预期增强，消费与投资韧性支撑增长。
  \item 全球通胀回落趋势确立，但核心服务通胀仍具粘性。
  \item 主要央行降息周期开启，但节奏和幅度存在不确定性。
\end{itemize}
\subsubsection*{分歧与条件差异}
\begin{itemize}
  \item 美联储与欧央行降息时点分歧：市场定价美联储9月首次降息，欧央行则可能更早行动。
  \item 对财政刺激效果的评估差异：部分机构认为美国财政扩张已见顶，另一些则强调基建和产业政策仍将提供支撑。
\end{itemize}
\subsubsection*{机构观点}
\paragraph{NOM} 创科实业受益于数据中心建设支出（1H26F同比增长22.1\%）、2020-21年销售高峰对应的5-7年替换周期启动以及渠道补库存三重动能。预计2026/27年每股盈利上调至0.76/0.87美元，毛利率从41.2\%提升至42.7\%。3Q26可能出现渠道补库存，因当前库存水平偏低且4Q终端销售有望超10\%增长。
{\small\color{refgray}数据：1H26F数据中心建设支出同比增长22.1\%\par}
{\small\color{refgray}数据：2026/27年每股盈利预测0.76/0.87美元\par}
{\small\color{refgray}数据：毛利率预计从41.2\%提升至42.7\%（2026年）\par}
{\small\color{refgray}数据：家得宝2026年6月电动工具库存同比持平或更低\par}
{\small\color{refgray}边际变化：从依赖家得宝单一渠道转向数据中心、欧洲等多元化增长，补库存预期从被动转为主动。\par}
\paragraph{GS} 埃克森美孚竞争优势从规模转向技术驱动的资产效率，轻质支撑剂和重新压裂技术正在提升Permian盆地单井采收率。管理层重申长期产量目标不变，但增长逻辑从'打更多井'转向'每口井出更多油'。下游炼化受益于低库存和韧性需求，全球成品油库存偏低对价格形成支撑。
{\small\color{refgray}数据：轻质支撑剂技术提升单井采收率约10\%\par}
{\small\color{refgray}数据：Permian盆地长期产量目标不变\par}
{\small\color{refgray}数据：全球成品油库存处于低位\par}
{\small\color{refgray}数据：资本配置优先级：强资产负债表、高回报增长、回购\par}
{\small\color{refgray}边际变化：市场关注点从产量增长转向技术驱动的资产效率提升，重新压裂技术延长井生命周期。\par}
\subsubsection*{关键数据}
\begin{itemize}
  \item 1H26F数据中心建设支出同比增长22.1\%
  \item 2026/27年每股盈利预测0.76/0.87美元（NOM）
  \item 毛利率预计从41.2\%提升至42.7\%（2026年）
  \item 轻质支撑剂技术提升单井采收率约10\%（GS）
  \item 全球成品油库存处于低位
\end{itemize}
{\small\color{refgray}本节综合 2 篇研究；涉及 NOM、GS。}
\newpage
\section{辅助信号｜战略咨询}
本板块聚焦波士顿咨询近期三份报告，提炼出印度制药创新拐点、工业数字化基础设施瓶颈、AI在能源交易中的差异化落地三大主题。
\subsection{印度制药：从仿制到原创的早期但非均衡转型}
\textbf{印度制药业正从低成本仿制向原创创新转型，但供应链依赖、数据基础设施薄弱和临床执行能力不足构成规模化瓶颈。}
\begin{itemize}
  \item 过去十年印度诞生10+个原创药物资产，专利提交量增长超4倍，生物技术初创从约1500家增至约2400家，PE/VC投资五年内增长约2.1倍至7.31亿美元。
  \item 政府投入约50亿美元，产学研联系加强，监管框架演进，共享研发与制造基础设施搭建，构成创新引擎雏形。
  \item 供应链对外依赖显著：基因治疗初创90\%材料从美欧中进口，交货周期30-45天；印度基因组计划仅1万个基因组，与英国生物银行约50万个差距巨大。
  \item 印度占全球约20\%疾病负担，但仅进行约4\%全球临床试验，临床执行能力不足与监管审批节奏不确定性使转化链条脆弱。
  \item 未来五年关键行动：培育专业生物技术资本、推动学术机构产品驱动、建立快速监管通道、建设本土供应链、弥合研发人才缺口。
\end{itemize}
\subsection{工业数字化：连接率仅13\%，基础设施先行}
\textbf{工业设备制造商数字化的真实瓶颈不在应用层，而在设备连接这一物理层基础，连接率仅13\%制约了上层应用的价值实现。}
\begin{itemize}
  \item 2020年调研显示，工业设备制造商与已安装设备基础的平均连接穿透率仅13\%，新设备连接率31\%，仍有35\%公司新设备未实现远程连接。
  \item 2016年企业优先预测性维护、远程监控等前台应用；2020年前十数字化优先级中四项变为使能技术：数据管理系统、智能设备、物联网骨干架构、设备连接性。
  \item 没有连接就没有数据流，预测性维护算法缺乏训练样本，远程诊断无从谈起；客户不愿分享运营数据，除非企业量化数据共享收益。
  \item 报告建议企业先解决设备在线问题，再部署上层应用，避免被AI预测、数字孪生等概念吸引而忽略基础。
\end{itemize}
\begin{figure}[htbp]
\centering
\includegraphics[width=0.88\linewidth]{figures/fig_185.jpg}
\caption{EXHIBIT 1 - Remote monitoring and diagnostics for field service technicians is also an attractive initial use of digital because of the nature of the business. Many industrial goods makers have large contingents of technicians who work on customer equipment. If companies }
\end{figure}
\begin{figure}[htbp]
\centering
\includegraphics[width=0.88\linewidth]{figures/fig_186.jpg}
\caption{EXHIBIT 2 - Some industrial goods companies offer incentives to help customers overcome their reluctance to connect new or previously installed equipment. One US heavy-machine manufacturer took a two-pronged approach to encourage customers to connect to its data platform,}
\end{figure}
\subsection{AI在能源交易：差异化落地，智能体层重塑运营}
\textbf{AI在能源交易中的价值取决于商品类型和工作流差异，智能体AI通过嵌入现有系统而非替代，实现从预测到执行的闭环价值。}
\begin{itemize}
  \item 充分利用AI的油气公司五年内可带来EBIT 30\%-70\%增量利润，但价值分布不均：电力和资金能源交易偏重预测性AI与优化，管道天然气、LNG和实物液体偏重智能体AI。
  \item 智能体AI架构在现有交易系统之上构建智能体层，可自动读取交易消息、提取条款、检查合规、准备录入、路由审批、起草确认函，减少多系统切换时间。
  \item 智能体AI的价值在于执行而非建议：更好的预测只有被交易台采纳才有价值，文档摘要只有加速审批或减少错误才有意义。
  \item 企业需明确规则：智能体自主完成什么、哪些只能准备、哪些必须人类审批，以及每一步如何记录保持问责。
\end{itemize}
\begin{figure}[htbp]
\centering
\includegraphics[width=0.88\linewidth]{figures/fig_189.jpg}
\caption{Exhibit 1 - \#\# The New Infrastructure Layer: Agents Operating Existing Tools To make this practical, firms need a trading IT architecture that connects AI to existing systems rather than bypassing them. Trading firms have plenty of tools. The problem is that people still }
\end{figure}
\begin{figure}[htbp]
\centering
\includegraphics[width=0.88\linewidth]{figures/fig_190.jpg}
\caption{Exhibit 2 - A recent BCG survey highlighted differences in adoption among trading organizations in power, pipeline gas, and physical liquids—the markets for which directly comparable benchmark data exists. (See Exhibit 2.) Power and pipeline gas traders lead in adoption, }
\end{figure}
\newpage
\section{辅助信号｜智库/国际机构}
本板块整合当日智库与国际机构最新研究，聚焦地缘政治风险定价、全球最低税实施效应、英国结构性财政约束、数字监管模型演进及贸易救济新动向，为宏观与产业决策提供辅助信号。
\subsection{地缘政治风险与主权信用定价的非对称性}
\textbf{威胁类地缘政治事件对新兴市场主权风险溢价的推升效应显著高于实际军事行动，市场对不确定性的定价比确定性事件更为激进。}
\begin{itemize}
  \item BIS工作论文基于2005-2025年13个新兴市场月度数据发现，威胁类地缘政治事件对主权CDS和EMBI利差的推升幅度比实际军事行动高出约30\%-50\%，且EMBI利差反应更为突出。
  \item 2022年俄罗斯入侵乌克兰后，新兴市场主权风险对地缘政治冲击的敏感度大幅上升，宏观基本面（债务水平、大宗商品贸易条件、全球资金条件）在特定地缘政治配置下被放大。
  \item 主权信用市场交易的是对未来违约概率的预期，威胁意味着不确定性上升，一旦威胁转化为行动，部分不确定性被消除，市场反而开始评估实际影响。
  \item 对组合管理有直接含义：当市场只关注已发生的地缘事件时，真正的不确定性溢价可能已体现在威胁阶段，需提前关注威胁类事件的定价信号。
\end{itemize}
\begin{figure}[htbp]
\centering
\includegraphics[width=0.88\linewidth]{figures/fig_161.jpg}
\caption{Figure 1 - Figure 1: Transmission of geopolitical risk shocks to sovereign risk premia in EMEs Nonetheless, given the complex nature of geopolitical shocks, several channels may affect the determinants of sovereign credit risk, and ex ante,}
\end{figure}
\begin{figure}[htbp]
\centering
\includegraphics[width=0.88\linewidth]{figures/fig_163.jpg}
\caption{Figure 3 - Figure 3: Response of SCDS spreads (left) and EMBI spreads (right) to a one-standard-deviation increase in the overall GPR index. Solid lines represent average responses evaluated at the panel-mean state. Shaded areas represent 68}
\end{figure}
\subsection{全球最低税实施首年：税率上升但投资与就业未现显著收缩}
\textbf{2024年全球最低税生效后，受覆盖跨国企业有效税率上升1.4-1.7个百分点，但集团层面的投资和就业未出现统计上显著的下降，政策冲击集中于高税筹倾向企业。}
\begin{itemize}
  \item OECD工作论文采用双重差分法，利用7.5亿欧元营收门槛两侧企业作为处理组和对照组，发现2024年实施期有效税率上升1.4-1.7个百分点，而2022-2023年公告期无显著提前调整行为。
  \item 税率上升并非均匀分布：此前低税率跨国企业有效税率上升约1.7个百分点，无形资产密集型企业同样是税率上升的集中领域，与GMT基于实质的carve-out设计逻辑一致。
  \item 论文未发现集团层面投资或就业的统计显著下降，这与事前模拟中“加税必然抑制投资”的线性推演形成对照，暗示企业可能通过其他渠道吸收税收成本。
  \item 当前结论需后续年份数据验证，但GMT不是一次性税收冲击，而是逐步嵌入企业决策框架的制度变量，未来低税利润规则（UTPR）全面生效后税收增量仍有上行空间。
\end{itemize}
\begin{figure}[htbp]
\centering
\includegraphics[width=0.88\linewidth]{figures/fig_174.jpg}
\caption{Figure 1 - \#\# 2.1 Overview of the GloBE rules The GMT applies to MNE groups with annual revenues of EUR 750m or more, with certain exclusions. Companies with more than EUR 750m in revenues in two of the four preceding years and with at least one foreign subsidiary are ge}
\end{figure}
\begin{figure}[htbp]
\centering
\includegraphics[width=0.88\linewidth]{figures/fig_176.jpg}
\caption{Figure 3 - \#\# 2.2 Timeline of GMT introduction The GMT was introduced as part of the discussions on a Two Pillar Solution to Address the Tax Challenges of the Digital Economy, under the auspices of the OECD/G20 Inclusive Framework on Base Erosion and Profit Shifting (BEP}
\end{figure}
\subsection{英国结构性约束：养老金财政不对称风险、能源价格错配与区域生产率差距}
\textbf{英国面临养老金三重锁制造的财政不对称风险、能源转型中的价格与投资错配、以及区域生产率差距固化三大结构性约束，财政空间有限要求政策排序与取舍。}
\begin{itemize}
  \item OECD经济调查通过随机模拟量化三重锁机制的财政后果：当工资或通胀意外上行时，养老金支出不成比例跳升，而经济放缓时几乎没有下调空间，基准情景可能系统性低估长期支出风险。
  \item 2025年英国一般政府债务占GDP比重102.3\%，利息支出持续上升，就业率偏低从收入端侵蚀养老金可持续性，自雇人群参与率不足加剧问题。
  \item 能源价格仍与国际天然气市场高度挂钩，高电价削弱电气化激励，延长对化石燃料价格波动的暴露，政策与网络投资错配是核心瓶颈。
  \item 区域生产率差距固化，报告建议通过地方财政自主权改革和技能培训体系落地来缩小差距，但财政空间有限要求优先排序：养老金改革、电网投资和分权化治理是未来十年关键叙事。
\end{itemize}
\begin{figure}[htbp]
\centering
\includegraphics[width=0.88\linewidth]{figures/fig_179.jpg}
\caption{Figure 1 - The United Kingdom faces a challenging economic environment, as past supply shocks and adverse external conditions, including high energy prices, trade tensions and the exit from the EU single market, continue to weigh on activity and public finances. The gove}
\end{figure}
\begin{figure}[htbp]
\centering
\includegraphics[width=0.88\linewidth]{figures/fig_180.jpg}
\caption{Figure 2 - Under the triple lock, state pensions have risen significantly faster than earnings, adding fiscal pressures and uncertainty, particularly during periods of macroeconomic volatility or weak growth (Figure 2). Reforming the triple lock is necessary to reduce fi}
\end{figure}
\subsection{数字监管模型谱系：从模式选择到条件匹配}
\textbf{各国数字监管正从单一模式转向混合设计，核心变量是公共与私人部门在目标设定与执行手段之间的责任分配，监管有效性取决于模型与条件的匹配度。}
\begin{itemize}
  \item OECD工作论文梳理六种监管模型（命令控制、绩效导向、市场机制、共同监管、元自律、行业自律），置于灵活性与公私责任坐标轴上，指出数字领域正形成“目标公管、手段私行”的主流实践。
  \item 监管模型选择应基于三类条件：对创新的影响、市场与私人部门特征、公共机构制度准备度。技术成熟度是常被低估的变量，对低成熟度技术施加刚性规则比不监管更危险。
  \item 该框架将监管选择从政策理念问题转化为条件匹配问题，产业决策者理解监管逻辑比预测监管内容更重要，因为监管逻辑决定了企业需要建立的合规能力和治理结构。
\end{itemize}
\begin{figure}[htbp]
\centering
\includegraphics[width=0.88\linewidth]{figures/fig_184.jpg}
\caption{Figure 1 - \#\# Spectrum of regulatory models Emphasising the respective role of public and private actors as a feature to classify models, while acknowledging that both these groups are always involved in the governance process in some capacity, is supported by the litera}
\end{figure}
\subsection{全球贸易救济新动向：过剩产能再分配与WTO谈判节奏}
\textbf{全球钢铁过剩产能通过贸易壁垒缝隙重新分配流向，南非等中等规模市场成为贸易摩擦新前线；WTO渔业补贴谈判进入关键阶段，数据缺口仍是最大障碍。}
\begin{itemize}
  \item 南非对进口冷轧钢发起保障措施调查，申请方安赛乐米塔尔南非公司（唯一本土生产商）指出全球产能过剩、主要地区加征壁垒、贸易流被迫转向更开放市场三因素叠加，导致进口激增。
  \item 该调查揭示可复用的观察框架：当多个主要市场同时收紧贸易壁垒时，过剩产能不会消失，只会流向壁垒最低的环节，南非正在积极使用WTO贸易救济工具。
  \item WTO渔业补贴谈判主席报告显示，核心议题已被充分分析，当前阻碍进展的是如何从已有方案中拼出多数成员能接受的组合；9月起将连续举办三轮“渔业周”，目标在2026年底前完成框架共识。
  \item 全球渔业补贴完整数据集仍然缺失，数据缺口不再是拖延借口，而是必须优先解决的短板，各成员按时提交准确通知是缩小信息差距的唯一途径。
\end{itemize}
\subsection{宏观政策与金融稳定}
\textbf{用于校准利率、通胀、财政约束和金融稳定叙事。}
\begin{itemize}
  \item 城市管理者面临一个共同困境：大量公共空间建成后迅速固化，无法回应社区不断变化的需求。经合组织最新区域发展报告提出一个反直觉的判断——真正提升城市韧性的，不是更多、更大的公共空间，而是那些能被社区集体塑造、管理和使用的“可适应公共空间”（Adaptive Communal Spac...
  \item 印度在节能领域已经走得很远。2025年，其一次能源强度改善超过4\%，是全球领先水平之一。到2026年3月，非化石能源装机容量已占53.21\%，提前接近2035年60\%的目标。但经合组织最新发布的《印度节能保险实施路线图》指出，真正制约节能研究从政策目标转化为市场行为的，不是技术，...
\end{itemize}
\begin{figure}[htbp]
\centering
\includegraphics[width=0.88\linewidth]{figures/fig_166.jpg}
\caption{Figure 3 - \#\# From tools to trajectory: a strategic approach The ACS policy framework is not a checklist in the narrow sense of predefined actions. It is a governance trajectory – a strategic learning process that connects civic initiative to institutional support throug}
\end{figure}
\begin{figure}[htbp]
\centering
\includegraphics[width=0.88\linewidth]{figures/fig_167.jpg}
\caption{Figure 3 - \textbackslash{}- MSE-SPICE (Scheme for Promotion and Investment in the Circular Economy): Provides a 25\% capital subsidy to MSMEs for adopting circular economy practices. \textbackslash{}- SIDBI financing scheme for energy savings projects: Targets MSMEs for energy efficiency projects wit}
\end{figure}
\subsection{AI、技术与生产率}
\textbf{用于判断 AI 投资周期、生产率扩散和产业约束是否正在改变资产定价。}
\begin{itemize}
  \item 全球公司治理正在经历一场静默但深刻的转向。经合组织最新发布的报告揭示了一个关键事实：截至2024年，60\%的司法管辖区允许公司发行不同投票权的股份，这一比例在2020年仅为44\%。更值得注意的是，92\%的辖区允许发行无投票权或有限投票权的优先股。这股潮流并非突然兴起，而是过去二十...
\end{itemize}
\begin{figure}[htbp]
\centering
\includegraphics[width=0.88\linewidth]{figures/fig_172.jpg}
\caption{Figure 1 - Evidence from the OECD Corporate Governance Factbook 2025 suggests that such flexibility is widespread: no less than \$90\textbackslash{}\%\$ of jurisdictions allow companies to issue shares with no voting rights (except for limited items). Ninety-two per cent also permit liste}
\end{figure}
\section{结语}
后续需验证的数据与叙事节点包括：美联储7月CPI数据对Warsh框架的检验、人民币中间价实际调整节奏、个股与指数波动率背离是否收敛、代理式AI对CPU需求的实质拉动、台积电与ASML毛利率能否持续高位、中国互联网公司利润修复的可持续性、全球纯电销量能否延续环比增长、Lp(a)靶点HORIZON试验结果、以及集运旺季需求能否支撑运价企稳。
\newpage
\section{覆盖概览}
投行/券商 87 篇；战略咨询 3 篇；智库/国际机构 9 篇。\par
投行覆盖：GS 25篇 / MS 22篇 / JPM 10篇 / NOM 10篇 / UBS 6篇 / Bernstein 4篇 / BofA 3篇 / Citi 3篇 / DB 2篇 / HSBC 1篇 / JEF 1篇\par
\section{Disclaimer}
{\scriptsize\color{refgray}
This community is a paid learning and information-sharing community created for general educational purposes only. The membership fee is charged solely for access to community discussions, educational content, organization, and operational costs. It is not an advisory fee, management fee, performance fee, brokerage fee, or compensation for personalized financial, investment, legal, tax, or professional advice.\par
I participate and share information solely as an individual and for educational discussion among community members. I am not acting as a financial adviser, investment adviser, broker, dealer, portfolio manager, legal adviser, tax adviser, accountant, fiduciary, or any other licensed professional adviser. Nothing shared in this community constitutes, and should not be interpreted as, financial advice, investment advice, legal advice, tax advice, a recommendation, solicitation, offer, endorsement, instruction, or guarantee to buy, sell, hold, short, trade, or invest in any security, cryptocurrency, fund, derivative, strategy, or other financial product.\par
All content shared in this community, including messages, discussions, links, files, charts, examples, market commentary, watchlists, AI-generated or AI-polished summaries, and third-party materials, is general, impersonal, and educational in nature. It does not take into account any individual's financial situation, investment objectives, risk tolerance, investment horizon, liquidity needs, tax status, legal restrictions, or suitability requirements. No content should be relied upon as suitable for any specific person, account, portfolio, or financial decision.\par
Information shared in this community may be incomplete, inaccurate, outdated, speculative, AI-generated, AI-edited, or based on third-party internet sources that have not been independently verified. I make no representation or warranty as to the accuracy, completeness, timeliness, reliability, or suitability of any information shared.\par
Financial markets involve substantial risk, including the possible loss of principal. Past performance, historical data, hypothetical examples, simulated results, backtests, personal opinions, or educational examples do not guarantee future results. Each member is solely responsible for conducting their own research, exercising independent judgment, and making their own decisions. Before making any financial, investment, legal, or tax decision, members should consult qualified and licensed professionals.\par
Participation in this community, payment of any membership fee, or communication with me or other members does not create any adviser-client, fiduciary, professional, contractual, agency, partnership, employment, or other advisory relationship. By joining or participating in this community, you acknowledge and agree that you are solely responsible for your own decisions, actions, transactions, profits, losses, risks, and consequences, and that I shall not be liable for any direct, indirect, incidental, consequential, financial, legal, tax, or other loss or damage arising from or related to any content, discussion, information, or material shared in this community.\par
}
\newpage
\begin{center}
{\Large 更多详情报告kcdesk.com}\par\vspace{0.8cm}
\includegraphics[width=0.58\linewidth]{\detokenize{../../prompts/zsxq_img.jpg}}
\end{center}
\end{document}
